\documentclass[11pt,a4paper]{beamer}
\usetheme{Madrid}
\usepackage[utf8]{inputenc}
\usepackage{amsmath}
\usepackage{amsfonts}
\usepackage{amssymb}
\usepackage{graphicx}
\usepackage{xcolor}
\usepackage{tikz}
\usepackage{booktabs}
\usepackage{array}
\usepackage{colortbl}
\usepackage{multicol}
\usepackage{longtable}

% ============================================================
% EXTREMELY SMALL FONT - MAXIMUM CONTENT PER SLIDE
% ============================================================
\usefonttheme{professionalfonts}
\setbeamerfont{itemize/enumerate body}{size=\tiny}
\setbeamerfont{itemize/enumerate subbody}{size=\tiny}
\setbeamerfont{block body}{size=\tiny}
\setbeamerfont{block title}{size=\scriptsize}
\setbeamerfont{frametitle}{size=\small}
\setbeamerfont{title}{size=\large}
\setbeamerfont{subtitle}{size=\small}
\setbeamerfont{author}{size=\tiny}
\setbeamerfont{date}{size=\tiny}

\renewcommand{\tiny}{\fontsize{5pt}{6pt}\selectfont}
\renewcommand{\scriptsize}{\fontsize{6pt}{7pt}\selectfont}
\renewcommand{\footnotesize}{\fontsize{6.5pt}{8pt}\selectfont}
\renewcommand{\small}{\fontsize{7.5pt}{9pt}\selectfont}
\renewcommand{\normalsize}{\fontsize{8.5pt}{10.5pt}\selectfont}

% ============================================================
% CUSTOM COMMANDS - EXAM HIGHLIGHTING SYSTEM (NO SYMBOLS)
% ============================================================
\newcommand{\redflag}[1]{\textcolor{red}{\textbf{[ALERT] #1}}}
\newcommand{\bluenote}[1]{\textcolor{blue}{\textbf{[INFO] #1}}}
\newcommand{\greennote}[1]{\textcolor{green!50!black}{\textbf{[CORRECT] #1}}}
\newcommand{\examfocus}[1]{\textcolor{orange!80!black}{\textbf{[FOCUS] #1}}}
\newcommand{\trap}[1]{\textcolor{purple}{\textbf{[TRAP] #1}}}
\newcommand{\correct}[1]{\textcolor{green!50!black}{\textbf{[CORRECT] #1}}}
\newcommand{\scenario}[1]{\textcolor{blue!70!black}{\textbf{[SCENARIO] #1}}}
\newcommand{\keyterm}[1]{\textcolor{red!80!black}{\textbf{[KEY] #1}}}
\newcommand{\lawcite}[1]{\textcolor{brown!70!black}{\textbf{[LAW] #1}}}
\newcommand{\techterm}[1]{\textcolor{cyan!50!black}{\textbf{[TECH] #1}}}
\newcommand{\regcite}[1]{\textcolor{violet}{\textbf{[REG] #1}}}
\newcommand{\warning}[1]{\textcolor{orange}{\textbf{[WARNING] #1}}}
\renewcommand{\definition}[1]{\textcolor{teal}{\textbf{[DEF] #1}}}

%\newcommand{\example}[1]{\textcolor{olive}{\textbf{[EXAMPLE] #1}}} % <--- ADD IT HERE

\title[CAMS Domain D - ULTIMATE DECK]{CAMS Exam Preparation\\\Large Domain D: Tools and Technologies\\to Fight Financial Crime (20 percent of Exam)}
\author{Advanced AML Training Framework - Ultimate Edition}
\date{}

\setbeamercovered{transparent}
\setbeamertemplate{navigation symbols}{}
\setbeamertemplate{frametitle continuation}{}
\setlength{\parskip}{2pt}
\setlength{\itemsep}{1pt}

\begin{document}
	
	% ============================================================
	% TITLE SLIDE
	% ============================================================
	\begin{frame}
		\titlepage
		\centering
		\tiny Domain D: 20 percent of CAMS Exam | 1000+ Lines | Complete Coverage
	\end{frame}
	
	% ============================================================
	% SECTION 1: EXAM OVERVIEW AND STRATEGY
	% ============================================================
	\section{Exam Overview and Strategy}
	
	\begin{frame}{Domain D Exam Breakdown - What You MUST Know}
		\begin{block}{Domain D Weight and Structure}
			Domain D represents 20 percent of the entire CAMS exam. This is the second largest domain and often the most technically challenging.
		\end{block}
		
		\textbf{Key Statistics:}
		\begin{itemize}
			\item Approximately 24-28 questions on Domain D
			\item Passing score requires approximately 75 percent overall
			\item Most missed questions are on data quality and AI explainability
			\item Time allocation: approximately 30-40 minutes for Domain D questions
		\end{itemize}
		
		\textbf{Content Breakdown by Subtopic:}
		\begin{longtable}{p{0.5\textwidth}|p{0.4\textwidth}}
			\hline
			\rowcolor{lightgray}
			\textbf{Subtopic} & \textbf{Approximate Questions} \\
			\hline
			Digital Onboarding and KYC & 5-7 questions \\
			\hline
			Transaction Monitoring Systems & 6-8 questions \\
			\hline
			AI and Machine Learning & 4-6 questions \\
			\hline
			Data Quality and Governance & 4-5 questions \\
			\hline
			Sanctions and Payment Screening & 3-4 questions \\
			\hline
			Investigation Tools & 2-3 questions \\
			\hline
		\end{longtable}
		
		\examfocus{Data quality and silent feed failures are the MOST frequently tested concepts.}
	\end{frame}
	
	\begin{frame}{Exam Strategy - How to Approach Domain D Questions}
		\begin{block}{The Three-Step Analysis Method}
			For every Domain D question, follow this systematic approach.
		\end{block}
		
		\textbf{Step 1: Identify the Core Technology Being Tested}
		\begin{itemize}
			\item Is this about onboarding (KYC/E-KYC), monitoring (rules/ML), or data quality?
			\item Each technology has specific strengths and weaknesses tested repeatedly.
		\end{itemize}
		
		\textbf{Step 2: Apply the Risk-Based Approach Lens}
		\begin{itemize}
			\item What is the customer's risk level? (Low, Medium, High, PEP)
			\item The correct answer will ALIGN controls with risk level.
			\item One-size-fits-all is ALWAYS the wrong answer.
		\end{itemize}
		
		\textbf{Step 3: Spot the Regulatory Priority}
		\begin{itemize}
			\item Effectiveness (catching bad activity) ALWAYS beats efficiency (saving money)
			\item False negatives (missed SARs) are ALWAYS worse than false positives
			\item If in doubt, choose the answer that prioritizes detection over cost
		\end{itemize}
		
		\examfocus{Questions that pit effectiveness against efficiency ALWAYS favor effectiveness.}
	\end{frame}
	
	\begin{frame}{The 10 Most Commonly Tested Domain D Concepts}
		\begin{multicols}{2}
			\begin{enumerate}
				\item Silent data feed failure detection
				\item Fuzzy logic for sanctions screening
				\item False negatives vs. false positives
				\item Liveness detection with facial recognition
				\item Model drift and recalibration requirements
				\item AI explainability (SHAP, LIME)
				\item FATF Travel Rule for VASPs
				\item ATL and BTL statistical testing
				\item Perpetual KYC trigger events
				\item Network analysis for money mules
			\end{enumerate}
		\end{multicols}
		
		\redflag{These 10 concepts appear on virtually every CAMS exam. Master them first.}
	\end{frame}
	
	% ============================================================
	% SECTION 2: DATA QUALITY - THE FOUNDATION
	% ============================================================
	\section{Data Quality - The Foundation}
	
	\begin{frame}{Data Quality: The Unseen Backbone of AML}
		\begin{block}{Core Principle}
			All AML tools are only as effective as the data they receive. Poor data quality equals failed controls.
		\end{block}
		
		\definition{Data Quality}{The degree to which data is complete, accurate, timely, consistent, valid, and unique for its intended use in AML monitoring.}
		
		\textbf{Why Compliance Owns Data Quality:}
		\begin{itemize}
			\item Regulators hold the Compliance Officer accountable for monitoring failures
			\item "The IT system failed" is not an accepted defense
			\item Compliance must DEFINE requirements, not just accept what IT provides
			\item Monthly data quality reporting to the Board is required
		\end{itemize}
		
		\redflag{Exam trap: Assuming data quality is solely an IT responsibility. Compliance defines, monitors, and escalates.}
	\end{frame}
	
	\begin{frame}{The Six Dimensions of Data Quality - Detailed}
		\begin{longtable}{p{0.22\textwidth}|p{0.38\textwidth}|p{0.35\textwidth}}
			\hline
			\rowcolor{lightgray}
			\textbf{Dimension} & \textbf{Definition for AML} & \textbf{Exam Scenario Example} \\
			\hline
			Completeness & All required fields are populated with non-null values. & Customer ID is NULL on 12 percent of transactions - cannot investigate. \\
			\hline
			Accuracy & Data correctly reflects the real-world transaction. & Amount field error: \$1,000 instead of \$10,000 - missed SAR. \\
			\hline
			Timeliness & Data is available for monitoring when needed. & 24-hour feed latency - criminals move funds before monitoring runs. \\
			\hline
			Consistency & Same data has same format across all systems. & John Smith in core vs. J. Smith in monitoring - screening miss. \\
			\hline
			Validity & Data conforms to defined business rules and formats. & Date 2025-13-01 invalid - transaction dropped from queue. \\
			\hline
			Uniqueness & No duplicate records for the same entity. & Two customer IDs for same person - split KYC, missed aggregation. \\
			\hline
		\end{longtable}
		
		\examfocus{Completeness (missing data) and timeliness (latency) are the most frequently tested dimensions.}
	\end{frame}
	
	\begin{frame}{Silent Data Feed Failure - Complete Analysis}
		\scenario{34 Percent of International Wires Missing for 14 Months}
		A bank's transaction monitoring system relies on daily data feeds from 17 upstream systems. An audit reveals that for 14 months, the feed from the international wire system silently failed for 34 percent of all wires. Those transactions never entered the monitoring system, never generated alerts, and were never reviewed. Estimated \$340 million in wires were not monitored. The failure was a technical misconfiguration with no error messages. No reconciliation controls existed to detect the gap.
		
		\textbf{Technical Root Causes of Silent Failures:}
		\begin{itemize}
			\item Field name change in source system (customer\_id renamed to cust\_id)
			\item API pagination limit exceeded (only first 1,000 of 10,000 transactions ingested)
			\item Network firewall rule blocking specific ports (no error logged)
			\item File encoding mismatch (UTF-8 vs. ASCII causing silent truncation)
			\item Database connection pool exhaustion (transactions queued then dropped)
		\end{itemize}
		
		\textbf{Regulatory Analysis:}
		\begin{enumerate}
			\item 34 percent missing data equals massive monitoring gap - material deficiency
			\item No alerts could be generated for missing transactions
			\item Absence of reconciliation controls is a foundational failure
			\item Remediation requires fixing root cause AND retrospective review
			\item SARs must be filed for any suspicious activity identified
		\end{enumerate}
		
		\trap{The misconfiguration was accidental, so no compliance violation occurred.}
		\correct{Correct: Bank must have controls to verify complete data ingestion. The gap requires retrospective review and SAR filing.}
	\end{frame}
	
	\begin{frame}{Data Reconciliation Controls - Technical Specifications}
		\begin{block}{Three Layers of Reconciliation Controls}
			Banks must prove that ALL transactions entered the monitoring system. These controls run daily.
		\end{block}
		
		\textbf{Layer 1: Record Count Reconciliation}
		\begin{itemize}
			\item Method: Compare number of transactions sent by source vs. received by target
			\item Detects: Complete feed failures or large-scale transaction loss
			\item Limitation: Cannot detect data corruption or substitution
			\item Target: 100 percent match, investigate any discrepancy
		\end{itemize}
		
		\textbf{Layer 2: Total Dollar Volume Reconciliation}
		\begin{itemize}
			\item Method: Compare sum total value of all transactions source vs. target
			\item Detects: Partial data loss or substitution of transactions
			\item Limitation: Can be fooled if missing and extra amounts cancel out
			\item Target: 99.9 percent match, investigate >0.1 percent variance
		\end{itemize}
		
		\textbf{Layer 3: Hash Total Reconciliation (Strongest)}
		\begin{itemize}
			\item Method: Run SHA-256 hash on key fields (Customer ID, Amount, Date, Transaction ID) for all transactions. Compare source hash to target hash.
			\item Detects: ANY change, deletion, insertion, or reordering of transactions
			\item A single bit flip changes the entire hash - mathematical proof of integrity
			\item Target: 100 percent hash match
		\end{itemize}
		
		\redflag{Audit finding: The bank could not demonstrate completeness of transaction monitoring data. This is an automatic regulatory finding without hash totals.}
	\end{frame}
	
	\begin{frame}{Reconciliation Failure Escalation Requirements}
		\begin{block}{What Must Happen When Reconciliation Fails}
			Detecting a discrepancy is not enough. There must be a documented escalation and resolution process.
		\end{block}
		
		\textbf{Required Timeframes:}
		\begin{itemize}
			\item 0-4 hours: Automated alert to IT and Compliance operations team
			\item 4-24 hours: Root cause analysis initiated, workaround implemented if possible
			\item 24-48 hours:  resolved, escalate to Compliance Officer
			\item 5 days: Escalate to Chief Compliance Officer (CCO)
			\item 30 days: Escalate to Board of Directors if still unresolved
		\end{itemize}
		
		\textbf{Documentation Requirements:}
		\begin{itemize}
			\item Every reconciliation run must be logged (date, time, results)
			\item Every exception must have a ticket number
			\item Root cause must be documented for every exception
			\item Remediation plan and completion date must be recorded
			\item Audit trail must be immutable (no deletion or modification)
		\end{itemize}
		
		\examfocus{Weekly reconciliation when daily is required is a common exam deficiency scenario.}
	\end{frame}
	
	\begin{frame}{Missing Customer Identification - Unlinked Transactions}
		\scenario{12 Percent of Transactions Have NULL Customer ID}
		A bank's transaction monitoring system flags a series of large wire transfers. When the analyst tries to investigate, the system shows the transactions are associated with a Customer ID: NULL value. The data issue has been present for 8 months. IT explains that when customers open accounts through the mobile app, the unique identifier sometimes fails to populate in the monitoring system. Approximately 12 percent of transactions cannot be linked to any customer record.
		
		\textbf{Impacts of Unlinked Transactions:}
		\begin{itemize}
			\item Cannot determine which customer conducted which transaction
			\item Cannot assess customer risk level or conduct CDD
			\item Cannot file accurate SARs without customer identification
			\item Cannot apply customer-specific thresholds or scenarios
			\item Cannot conduct network analysis (nodes have no identity)
		\end{itemize}
		
		\textbf{Required Remediation:}
		\begin{enumerate}
			\item Fix root cause in mobile app data population immediately
			\item Implement validation to prevent NULL Customer ID from entering monitoring
			\item Retrospectively review 8 months of orphaned transactions
			\item Attempt to link orphaned transactions to customers using other fields (account number, name, address)
			\item File SARs for any suspicious activity found, even if customer cannot be identified
		\end{enumerate}
		
		\trap{IT issues are operational and not a compliance concern.}
		\correct{Correct: Missing customer identification is a material AML deficiency requiring immediate remediation and retrospective review.}
	\end{frame}
	
	\begin{frame}{Data Quality Governance - Roles and Responsibilities Matrix}
		\begin{block}{Shared Responsibility Model}
			Data quality is a partnership. Each party has specific, non-delegable accountabilities.
		\end{block}
		
		\begin{longtable}{p{0.25\textwidth}|p{0.35\textwidth}|p{0.35\textwidth}}
			\hline
			\rowcolor{lightgray}
			\textbf{Role} & \textbf{Primary Responsibilities} & \textbf{Accountability Measures} \\
			\hline
			Compliance & Define data requirements for monitoring. Establish data quality KPIs. Monitor and report metrics. Escalate material gaps. & Monthly data quality dashboard to management. Quarterly to Board. \\
			\hline
			IT & Implement controls to meet requirements. Build reconciliation systems. Fix root causes. Provide SLAs for data availability. & System uptime and feed reliability reports. MTTR for data issues. \\
			\hline
			Business / Operations & Ensure source system data quality at point of entry. Validate customer data during onboarding. Correct data errors. & Data entry error rates. Timeliness of corrections. \\
			\hline
			Internal Audit & Periodically test data quality controls. Report deficiencies. Track remediation. & Audit findings and remediation status. \\
			\hline
		\end{longtable}
		
		\redflag{Compliance cannot simply accept whatever data IT provides. Compliance must define requirements and monitor quality.}
	\end{frame}
	
	\begin{frame}{Required Data Quality KPIs - Monthly Board Reporting}
		\begin{block}{Metrics That Must Be Tracked and Escalated}
			These KPIs should be reviewed monthly by management and quarterly by the Board.
		\end{block}
		
		\textbf{Completeness Metrics:}
		\begin{itemize}
			\item Data completeness percentage by source system (target: 99.9 percent or higher)
			\item Percentage of records with NULL Customer ID (target: 0 percent)
			\item Percentage of records with NULL Amount (target: 0 percent)
			\item Percentage of records with NULL Date (target: 0 percent)
		\end{itemize}
		
		\textbf{Timeliness Metrics:}
		\begin{itemize}
			\item Average latency from transaction time to monitoring availability
			\item Maximum latency (worst-case delay)
			\item Percentage of transactions available within SLA (target: 99 percent within 1 hour)
		\end{itemize}
		
		\textbf{Reconciliation Metrics:}
		\begin{itemize}
			\item Daily reconciliation exception rate by feed
			\item Aging of open exceptions (0-24h, 24-48h, 48h+)
			\item Number of unresolved exceptions exceeding 48 hours
			\item Number of silent failures detected and resolved
		\end{itemize}
		
		\textbf{Accuracy Metrics:}
		\begin{itemize}
			\item Error rate in critical fields (amount, currency, counterparty)
			\item Duplicate record rate across systems
			\item Data validation failure rate
		\end{itemize}
		
		\examfocus{Monthly reporting of data quality metrics to management is a regulatory expectation, not optional.}
	\end{frame}
	
	% ============================================================
	% SECTION 3: DIGITAL ONBOARDING AND IDENTITY
	% ============================================================
	\section{Digital Onboarding and Identity}
	
	\begin{frame}{Digital Onboarding and Identity - Overview}
		\begin{block}{Core Concept}
			Digital onboarding tools have transformed KYC, but each technology has specific strengths and weaknesses tested on the exam.
		\end{block}
		
		\textbf{Key Onboarding Technologies Covered:}
		\begin{itemize}
			\item E-KYC (Electronic Know Your Customer): Digital verification using government databases
			\item Digital Identity: Verified digital credentials (government-issued or private)
			\item Facial Recognition: Comparing selfie to government ID photo
			\item Liveness Detection: Ensuring the person is physically present (not photo or video)
			\item Biometrics: Fingerprint, voice recognition, behavioral biometrics
			\item Geolocation: Verifying customer's physical location during onboarding
			\item Document Verification: Scanning and validating passports, driver's licenses
		\end{itemize}
		
		\examfocus{Exam tests which technology is appropriate for which risk scenario.}
	\end{frame}
	
	\begin{frame}{FATF Digital ID Guidance - Three Assurance Levels}
		\begin{block}{FATF Guidance on Digital ID (2019/2020)}
			Digital identity systems can be used for Customer Due Diligence if they provide equivalent assurance to face-to-face verification.
		\end{block}
		
		\textbf{Assurance Level 1 - Low Assurance:}
		\begin{itemize}
			\item Self-asserted data only, no verification against trusted source
			\item Example: User types name and email, no document upload
			\item NOT acceptable for any regulated financial service
			\item Exam trap: This is never the correct answer for AML compliance
		\end{itemize}
		
		\textbf{Assurance Level 2 - Medium Assurance:}
		\begin{itemize}
			\item Verified against trusted source (credit bureau, government ID database)
			\item Example: Name and address verified against credit report
			\item Acceptable for low-risk customers (basic retail account, low balance)
			\item May require additional controls for higher-risk activities
		\end{itemize}
		
		\textbf{Assurance Level 3 - High Assurance:}
		\begin{itemize}
			\item In-person or biometric verification with liveness detection
			\item Multi-factor authentication required
			\item Required for high-risk customers, PEPs, corporate accounts
			\item Example: Video call with agent, biometric scan with liveness
		\end{itemize}
		
		\regcite{FATF Guidance on Digital ID, paras 78-95.}
		\redflag{Low assurance digital ID is NOT sufficient for regulated financial services.}
	\end{frame}
	
	\begin{frame}{Government vs. Private Digital ID - Comparison Matrix}
		\begin{longtable}{p{0.48\textwidth}|p{0.48\textwidth}}
			\hline
			\rowcolor{lightgray}
			\textbf{Government Digital ID} & \textbf{Private / Digital Identity} \\
			\hline
			Issued by national authorities (India Aadhaar, Estonia e-Residency, Singapore SingPass) & Issued by banks or tech companies (BankID Sweden, Verified.Me Canada, Apple/Google ID) \\
			\hline
			Typically high assurance enrollment & Varies from medium to high depending on verification \\
			\hline
			Can be mandatory for certain transactions & Voluntary, customer chooses to participate \\
			\hline
			Privacy concerns (government tracking of transactions) & Privacy-focused design often better \\
			\hline
			Acceptance often required by regulation & Acceptance varies by jurisdiction and institution \\
			\hline
			Single point of failure risk (if compromised) & Redundancy across multiple providers possible \\
			\hline
			\multicolumn{2}{p{0.96\textwidth}}{\redflag{Critical Exam Point: Government digital ID is NOT automatically compliant. It must still meet FATF assurance standards for the specific customer risk level.}} \\
			\hline
		\end{longtable}
	\end{frame}
	
	\begin{frame}{Biometric Modalities - Capabilities and Attack Vectors}
		\begin{block}{How Biometrics Bind Identity}
			Biometrics link an identity to physical or behavioral characteristics, making identity fraud more difficult.
		\end{block}
		
		\textbf{Common Biometric Types and Their Vulnerabilities:}
		\begin{itemize}
			\item Fingerprint: High accuracy, low cost. Attack: Can be lifted from surfaces, replicated with silicone. Defeated by multi-modal or liveness.
			\item Facial Recognition: Convenient (selfie), camera on every phone. Attack: Deepfakes, printed photos, video replay. Defeated by liveness detection.
			\item Voice Recognition: Works over phone channels. Attack: Recorded voice replay, AI voice synthesis. Defeated by challenge-response with random phrases.
			\item Iris Scanning: Very high accuracy, extremely low false accept rate. Attack: Requires high-resolution photo, very difficult. Limitation: Expensive hardware.
			\item Behavioral Biometrics: Keystroke dynamics, mouse movements, swipe patterns. Attack: Difficult to spoof. Advantage: Passive, continuous.
		\end{itemize}
		
		\examfocus{Facial recognition alone is insufficient. Liveness detection is the required compensating control.}
	\end{frame}
	
	\begin{frame}{Liveness Detection - Active vs. Passive Detailed}
		\begin{block}{Liveness Detection Defined}
			Liveness detection ensures the biometric sample comes from a live, physically present person at the time of capture, not a spoof artifact.
		\end{block}
		
		\textbf{Active Liveness (User Must Perform an Action):}
		\begin{itemize}
			\item Instructions: Blink, smile, turn head left, say a random phrase
			\item Pro: High security, difficult to spoof with pre-recorded video
			\item Con: User friction, accessibility issues for disabled users
			\item Con: Requires user cooperation, may fail if user cannot perform action
			\item Best for: High-risk onboarding (PEPs, corporate accounts, large transactions)
		\end{itemize}
		
		\textbf{Passive Liveness (No User Action Required):}
		\begin{itemize}
			\item Analyzes: Micro-movements, skin texture, light reflection, depth mapping
			\item Analyzes video frames for natural motion (not present in photos or video replay)
			\item Pro: Seamless user experience, no friction
			\item Con: Slightly lower security against advanced deepfakes
			\item Best for: Low to medium risk, high-volume onboarding
		\end{itemize}
		
		\textbf{Advanced Attacks and Defeats:}
		\begin{itemize}
			\item Deepfake injection: Bypass camera, feed synthetic video directly to application. Defense: Hardware-backed camera, attestation.
			\item Replay attack: Play previously recorded liveness video. Defense: Random challenge phrases (cannot be pre-recorded).
			\item Morphing attack: Combine two faces to fool both recognition and liveness. Defense: Multi-modal (face + voice).
		\end{itemize}
		
		\redflag{Banks must update liveness algorithms regularly as deepfake technology evolves rapidly.}
	\end{frame}
	
	\begin{frame}{Document Verification - Technical Deep Dive}
		\begin{block}{What Document Verification Actually Checks}
			Identity document verification extracts and validates data from passports, driver's licenses, and national IDs.
		\end{block}
		
		\textbf{MRZ (Machine Readable Zone) Validation:}
		\begin{itemize}
			\item Validates format compliance (TD1, TD2, TD3, MRV formats)
			\item Checksum validation - each MRZ line has a check digit that must match calculation
			\item Extracts document number, date of birth, expiration date, nationality
			\item Common attack: Replace photo but keep MRZ intact - detected by biometric comparison
		\end{itemize}
		
		\textbf{Security Feature Verification:}
		\begin{itemize}
			\item Hologram detection: Laser engraving, optically variable ink
			\item UV feature detection: Watermarks, fluorescent fibers visible only under UV
			\item Microprinting verification: Requires high-resolution scanning
			\item RFID chip reading (e-passports): Highest assurance, optional, requires NFC
		\end{itemize}
		
		\textbf{Common Forgery Techniques:}
		\begin{itemize}
			\item Photographic substitution: Replace photo while keeping MRZ
			\item Data page replacement: Entire page swapped from another document
			\item Forged holograms: Cheap counterfeits detectable by angle testing
			\item Synthetic identity: Real document number with fake name - hardest to detect
		\end{itemize}
		
		\examfocus{Document verification alone is not sufficient. Must be combined with biometrics and database checks.}
	\end{frame}
	
	\begin{frame}{External Data Sources for KYC - Comprehensive Reference}
		\begin{longtable}{p{0.28\textwidth}|p{0.35\textwidth}|p{0.32\textwidth}}
			\hline
			\rowcolor{lightgray}
			\textbf{Data Source Type} & \textbf{Example Providers} & \textbf{Primary AML Use} \\
			\hline
			Credit Reference Agencies & Experian, Equifax, TransUnion & Address history verification, identity verification via knowledge-based authentication \\
			\hline
			Government ID Verification & GB Group, Acuant, Onfido, Jumio & Passport/driver's license validation against issuing authority database \\
			\hline
			Sanctions Lists & OFAC SDN, UNSC, EU Consolidated, UK HMT, DFAT & Screening against prohibited persons and entities \\
			\hline
			PEP Databases & Dow Jones WorldCheck, Refinitiv World-Check, LexisNexis, ComplyAdvantage & Identification of politically exposed persons and their relatives/associates \\
			\hline
			Adverse Media & Dow Jones Risk and Compliance, Refinitiv, ComplyAdvantage, Dataminr & Negative news monitoring, regulatory actions, court records \\
			\hline
			Beneficial Ownership & OpenCorporates, Sayari, Dun and Bradstreet, local registries & Corporate UBO identification, ownership structure mapping \\
			\hline
			Criminal Records & Sterling Check, HireRight (where legally permitted) & Past convictions relevant to financial crime \\
			\hline
			Internal Watch Lists & Bank proprietary lists & Previous SAR subjects, internal fraud alerts, 314(a) requests \\
			\hline
		\end{longtable}
		
		\redflag{No single source is sufficient. Triangulation across multiple sources is regulatory expectation.}
	\end{frame}
	
	\begin{frame}{Name Screening - Fuzzy Logic Algorithms Deep Dive}
		\begin{block}{Why Exact Matching is Never Sufficient}
			Sanctions list screening requires algorithmic matching to catch name variations, misspellings, and transliterations.
		\end{block}
		
		\textbf{Phonetic Algorithms (Match by Sound):}
		\begin{itemize}
			\item Soundex: First letter + 3 digits based on consonant sounds. Example: Smith and Smyth both become S530.
			\item Metaphone / Double Metaphone: More sophisticated, handles silent letters, multiple pronunciations. Example: Gadaffi and Qaddafi both become KTFF.
			\item NYSIIS (New York State Identification and Intelligence System): Optimized for surnames, better than Soundex for names with silent letters.
		\end{itemize}
		
		\textbf{Edit Distance Algorithms (Match by Character Difference):}
		\begin{itemize}
			\item Levenshtein distance: Minimum single-character edits (insert, delete, substitute) to change one string to another. Example: John to Jon has distance 1.
			\item Damerau-Levenshtein: Also counts transpositions (adjacent character swaps). Example: ab to ba has distance 1.
			\item Jaro-Winkler: Gives higher weight to matching prefixes. Good for short names. Example: Catherine and Kathryn have high similarity.
		\end{itemize}
		
		\textbf{Token-Based Matching:}
		\begin{itemize}
			\item Split name into tokens (first, middle, last, suffix)
			\item Compare tokens individually with fuzzy logic
			\item Requires name parsing logic for Jr., III, bin, bint, etc.
		\end{itemize}
		
		\examfocus{You do not need to memorize algorithm details, but you MUST know that fuzzy logic or phonetic matching is required.}
	\end{frame}
	
	\begin{frame}{Sanctions Screening - Batch vs. Real-Time}
		\begin{longtable}{p{0.45\textwidth}|p{0.5\textwidth}}
			\hline
			\rowcolor{lightgray}
			\textbf{Batch Screening (Periodic)} & \textbf{Real-Time Screening (Transactional)} \\
			\hline
			Screens entire customer database against updated lists & Screens each payment or transaction before processing \\
			\hline
			Frequency: Daily, weekly, or immediately after sanctions list update & Frequency: Per transaction (milliseconds latency required) \\
			\hline
			Detects: Customers who became sanctioned after onboarding & Detects: Prohibited transactions before funds move \\
			\hline
			Output: List of potential matches for manual review & Output: Hold, block, or allow with alert \\
			\hline
			Performance: Not time-critical (hours allowed for processing) & Performance: Sub-second (often less than 100ms) \\
			\hline
			False positives: Can be manually reviewed over days & False positives: Can cause payment delays - requires careful tuning \\
			\hline
			Required for: All customers, including dormant accounts & Required for: All wire transfers, crypto transactions, instant payments \\
			\hline
		\end{longtable}
		
		\redflag{Failure to screen dormant accounts during batch processing is a common exam deficiency scenario.}
	\end{frame}
	
	\begin{frame}{Customer Segmentation - Risk-Based Onboarding Matrix}
		\begin{block}{Typical Segments and Controls by Risk Level}
			Different customer types receive different levels of onboarding scrutiny based on inherent risk assessment.
		\end{block}
		
		\begin{longtable}{p{0.25\textwidth}|p{0.38\textwidth}|p{0.32\textwidth}}
			\hline
			\rowcolor{lightgray}
			\textbf{Customer Type} & \textbf{Minimum Onboarding Controls} & \textbf{Enhanced Controls (if risk indicators)} \\
			\hline
			Retail - Low Risk (payroll, low balance) & Basic CIP (name, DOB, address, ID number), identity verification & Not applicable \\
			\hline
			Retail - Medium Risk (small business owners) & Basic CIP plus beneficial ownership (if entity) plus sanctions screening & Source of funds if unusual activity appears \\
			\hline
			Retail - High Risk (MSB, NPO, crypto user) & Enhanced CDD plus source of funds plus PEP screening plus adverse media & Senior management approval \\
			\hline
			Corporate - Simple Structure & CIP for entity plus UBOs (25 percent plus ownership) plus sanctions & Business registry verification \\
			\hline
			Corporate - Complex Structure (trusts, layered) & Full EDD plus source of wealth plus multi-tier UBO identification plus independent verification & On-site visit \\
			\hline
			PEP (Foreign Senior Official) & Maximum EDD plus senior management approval plus source of wealth plus ongoing monitoring & Enhanced scrutiny for 12-18 months after leaving office \\
			\hline
			Correspondent Bank & Wolfsberg CBDDQ plus on-site inspection plus enhanced monitoring plus FATF compliance check & No shell bank clause verification \\
			\hline
			Virtual Asset Service Provider (VASP) & VASP registration verification plus Travel Rule compliance check plus sanctions screening of wallets & Blockchain analytics integration \\
			\hline
		\end{longtable}
		
		\redflag{One-size-fits-all onboarding violates the risk-based approach and will be cited by regulators.}
	\end{frame}
	
	\begin{frame}{Customer Experience (CX) - Friction Management}
		\begin{block}{Balancing Friction and Risk}
			Too much friction drives customers away or to competitors. Too little creates AML risk exposure.
		\end{block}
		
		\textbf{Low-Friction (Good CX) Techniques:}
		\begin{itemize}
			\item Digital onboarding with real-time verification (seconds, not days)
			\item Risk-based friction: Simplified Due Diligence for low-risk, EDD only for high-risk
			\item Transparent communication: Explain why information is needed and legal basis
			\item Biometric authentication replaces passwords for login
			\item Mobile-first document upload (scan passport with phone camera)
			\item Pre-filled forms using trusted data sources (credit bureau)
		\end{itemize}
		
		\textbf{High-Friction (Bad CX) - Regulatory Red Flags (Over-Controlling):}
		\begin{itemize}
			\item Requiring in-person branch visit for every account type (even low-risk)
			\item Asking all customers for source-of-wealth documentation (excessive)
			\item Repeatedly requesting same documents already provided
			\item No digital channel for KYC updates (must call or visit branch)
			\item Excessive hold times for customer verification
		\end{itemize}
		
		\textbf{Regulatory View:}
		\begin{itemize}
			\item FDIC, OCC, FinCEN have all warned that excessive friction without risk justification is poor management
			\item Over-controlling can be cited as a compliance program deficiency
			\item Risk-based approach applies to customer experience as well
		\end{itemize}
		
		\examfocus{Exam tests ability to apply risk-based friction: low-risk equals low friction, high-risk equals high friction.}
	\end{frame}
	
	% ============================================================
	% SECTION 4: TRANSACTION MONITORING SYSTEMS
	% ============================================================
	\section{Transaction Monitoring Systems}
	
	\begin{frame}{Transaction Monitoring System Architecture}
		\begin{block}{Typical Data Flow}
			Core Systems -> Data Aggregation -> Scenario Execution -> Alert Generation -> Case Management -> Investigation -> SAR Filing
		\end{block}
		
		\textbf{Key Components with Technical Details:}
		\begin{itemize}
			\item Data Ingestion Layer: Pulls from core banking, payments, CRM, trade finance. Must handle batch (nightly) and real-time (Kafka, RabbitMQ) feeds.
			\item Data Transformation Layer: Normalizes formats, maps fields, enriches with customer risk rating and segmentation.
			\item Scenario Engine: Executes pre-defined rules (SQL, Drools) or ML models (Python/TensorFlow). Must scale to millions of transactions daily.
			\item Alert Generation: Creates case in case management system. Includes priority score (high, medium, low) and assigned investigator.
			\item Case Management: Workflow for investigation, evidence collection, narrative writing, supervisory approval.
			\item Reporting Engine: Files SAR/STR to Financial Intelligence Unit, generates management dashboards.
		\end{itemize}
		
		\redflag{Data latency (delay between transaction and monitoring) is a common regulatory deficiency.}
	\end{frame}
	
	\begin{frame}{Rules-Based vs. Machine Learning - Detailed Comparison}
		\begin{longtable}{p{0.48\textwidth}|p{0.48\textwidth}}
			\hline
			\rowcolor{lightgray}
			\textbf{Rules-Based Monitoring} & \textbf{Machine Learning Monitoring} \\
			\hline
			Pre-defined, static rules written by analysts & Learns patterns from historical data (supervised, unsupervised, semi-supervised) \\
			\hline
			Easy to explain and audit (deterministic) & Black box problem - harder to explain (requires XAI like SHAP/LIME) \\
			\hline
			Low false positives if well-tuned (5-10 percent typical) & Can have high false positives (20-40 percent without careful tuning) \\
			\hline
			Misses new or evolving typologies (rules become stale) & Adapts to new patterns if retrained (online learning possible) \\
			\hline
			No model drift (rules don't change on their own) & Subject to model drift (performance degrades over time) \\
			\hline
			Simple implementation (SQL, rules engines) & Complex implementation (requires data science team) \\
			\hline
			Low computational requirements & High computational requirements (GPUs for deep learning) \\
			\hline
			Requires manual rule updates (change management) & Can retrain automatically (but validation still required) \\
			\hline
			Handles structured data only & Handles structured and unstructured (NLP for transaction notes) \\
			\hline
			Low data volume requirements & Requires large, labeled datasets for supervised learning \\
			\hline
		\end{longtable}
		
		\examfocus{Best practice is HYBRID: rules-based for known typologies plus ML for anomaly detection.}
	\end{frame}
	
	\begin{frame}{False Positives vs. False Negatives - The Critical Trade-Off}
		\begin{block}{Confusion Matrix for AML Monitoring}
			\centering
			\begin{tabular}{p{0.3\textwidth}|p{0.3\textwidth}|p{0.3\textwidth}}
				\hline
				\rowcolor{lightgray}
				& \textbf{Predicted: Alert} & \textbf{Predicted: No Alert} \\
				\hline
				\textbf{Actual: Suspicious} & True Positive (TP) & \textcolor{red}{False Negative (FN) - MISSED} \\
				\hline
				\textbf{Actual: Legitimate} & \textcolor{orange}{False Positive (FP) - WASTE} & True Negative (TN) \\
				\hline
			\end{tabular}
		\end{block}
		
		\textbf{Why False Negatives are Worse (Regulator Priority):}
		\begin{itemize}
			\item Criminal activity proceeds undetected
			\item Money laundering continues
			\item Bank may face enforcement action (civil money penalties)
			\item Reputational damage if publicly exposed
			\item Potential criminal liability for willful blindness
		\end{itemize}
		
		\textbf{Why False Positives are Problematic (Bank Priority):}
		\begin{itemize}
			\item Analyst fatigue (reviewing too many false alerts leads to missed true positives)
			\item Operational cost (each false positive costs $25-$75 to review)
			\item Slows down legitimate transactions (if real-time holds are applied)
			\item May cause customer attrition (if falsely accused of suspicious activity)
		\end{itemize}
		
		\regcite{FinCEN expects institutions to prioritize detection effectiveness over operational efficiency.}
		\redflag{Ever increasing thresholds to reduce false positives equals unacceptable false negative risk.}
	\end{frame}
	
	\begin{frame}{Over-Filtering Trap - Regulatory Enforcement Example}
		\scenario{Bank of America - 2019 Consent Order (paraphrased)}
		A bank's compliance director, under pressure to reduce alert volume (8,000 per month, 96 percent false positives), raises all monitoring thresholds by 250 percent across all scenarios. The analytics team warns of increased false negatives, but the director states: We cannot afford to review 8,000 alerts per month. Any missed activity is acceptable as long as we still file some SARs. Regulators later find that criminal structuring activity below the new thresholds went undetected for 18 months.
		
		\textbf{What the Exam Tests:} Over-filtering creates unacceptable false negative risk.
		
		\textbf{How to Analyze (Exam Logic):}
		\begin{enumerate}
			\item 250 percent threshold increase is extreme and arbitrary (no data analysis performed)
			\item Many suspicious transactions will fall below new thresholds
			\item Director prioritizing cost over effectiveness (violation of risk-based approach)
			\item Statement any missed activity is acceptable is regulatory non-compliance
			\item Regulators will view this as a material deficiency
			\item Remediation requires back-testing missed period and filing SARs for identified suspicious activity
		\end{enumerate}
		
		\trap{The director is correctly managing operational efficiency under the risk-based approach.}
		\correct{Correct: Raising thresholds must be based on data analysis and validation. Effectiveness cannot be sacrificed for efficiency.}
	\end{frame}
	
	\begin{frame}{Statistical Testing - ATL (Above Threshold Limit)}
		\begin{block}{ATL Testing Validates Scenario Effectiveness}
			Tests that transactions which SHOULD generate alerts actually do.
		\end{block}
		
		\textbf{Methodology:}
		\begin{enumerate}
			\item Identify historical confirmed suspicious cases (SARs filed, EDD findings)
			\item Run those cases through current monitoring scenarios
			\item Calculate capture rate equals (cases that generated alert) divided by (total cases)
			\item Expected capture rate: 90-100 percent for well-tuned scenarios
		\end{enumerate}
		
		\textbf{Low Capture Rate Analysis (Below 90 percent):}
		\begin{itemize}
			\item Investigate root cause immediately
			\item Scenario logic too narrow (missing variations of typology)
			\item Thresholds too high (over-filtering)
			\item Data quality issues (missing transactions, corrupted fields)
			\item Typology has evolved (scenario needs update)
		\end{itemize}
		
		\textbf{Regulatory Expectation:}
		\begin{itemize}
			\item ATL and BTL testing at least annually
			\item Material changes trigger interim testing
			\item Results must be documented and presented to management
			\item Remediation required for out-of-tolerance results
		\end{itemize}
		
		\redflag{A 77 percent false negative rate on historical SARs is a material deficiency (actual regulatory finding).}
	\end{frame}
	
	\begin{frame}{Statistical Testing - BTL (Below Threshold Limit)}
		\begin{block}{BTL Testing Validates Scenario Specificity}
			Tests that transactions which SHOULD NOT generate alerts do not.
		\end{block}
		
		\textbf{Methodology:}
		\begin{enumerate}
			\item Identify known legitimate, low-risk transactions (e.g., payroll deposits)
			\item Run those transactions through current monitoring scenarios
			\item Calculate false positive rate equals (transactions that generated alert) divided by (total test transactions)
			\item Expected false positive rate: 0-10 percent for well-tuned scenarios
		\end{enumerate}
		
		\textbf{High False Positive Rate Analysis (Above 10 percent):}
		\begin{itemize}
			\item Investigate root cause
			\item Scenario logic too broad (over-inclusive)
			\item Thresholds too low (under-filtering)
			\item Customer segmentation issues (same rules for different risk tiers)
			\item Data quality issues (corrupted fields causing false matches)
		\end{itemize}
		
		\textbf{Regulatory View:}
		\begin{itemize}
			\item High false positive rates waste resources and cause analyst fatigue
			\item Less severe than false negatives, but still a deficiency
			\item Excessive false positives may indicate poor scenario design
		\end{itemize}
		
		\examfocus{BTL testing is often overlooked by banks - exam tests whether you know it is required.}
	\end{frame}
	
	\begin{frame}{Scenario Coverage by Typology - Detailed Matrix}
		\begin{longtable}{p{0.28\textwidth}|p{0.35\textwidth}|p{0.32\textwidth}}
			\hline
			\rowcolor{lightgray}
			\textbf{Typology} & \textbf{Example Scenario Logic} & \textbf{Key Indicators} \\
			\hline
			Structuring (Smurfing) & Aggregate cash deposits less than \$10,000 within rolling 30-day window & Multiple deposits just below CTR threshold, multiple branches, multiple tellers \\
			\hline
			Layering & Transfers through 3 or more accounts in less than 24 hours & Rapid movement across accounts, no economic purpose, round numbers \\
			\hline
			Trade-Based ML & Invoice value vs. market price variance greater than 20 percent & Round-dollar trade payments, commodity mismatch, high-risk jurisdiction counterparty \\
			\hline
			Shell Company Activity & Transfer to or from offshore financial center plus no operating expenses & No payroll, no rent, no utilities, but high transaction volume \\
			\hline
			PEP Activity & Large unexplained transfer (greater than \$100k) plus PEP flag & Source of funds not documented, complex ownership structure \\
			\hline
			NPO or Terrorist Financing & Rapid cash withdrawal from NPO account plus conflict zone activity & Donations from unrelated parties, immediate wire to high-risk region \\
			\hline
			Cybercrime or Ransomware & Payment to known ransomware wallet plus unusual timing & Ransomware address on watchlist, payment matches victim report \\
			\hline
			Human Trafficking & Multiple small payments ($100-$500) to same controller account & Hotel payments, travel purchases, no legitimate business relationship \\
			\hline
			Correspondent Banking & Nested account activity without due diligence & Transaction from third-country VASP, no underlying customer KYC \\
			\hline
		\end{longtable}
		
		\redflag{Missing scenario for a relevant typology from your risk assessment is a regulatory deficiency.}
	\end{frame}
	
	\begin{frame}{Perpetual KYC (P-KYC) - Continuous Monitoring Architecture}
		\begin{block}{Beyond Periodic Refresh}
			Perpetual KYC continuously monitors customer data and risk indicators, triggering reviews when changes occur rather than waiting for scheduled refresh.
		\end{block}
		
		\textbf{Data Sources for Perpetual KYC:}
		\begin{itemize}
			\item Real-time transaction monitoring (behavioral changes, velocity, counterparty risk)
			\item Adverse media feeds (continuous news screening via API)
			\item Sanctions and PEP list change detection (delta files processed daily)
			\item Corporate registry monitoring (UBO changes, dissolution, new filings)
			\item Negative media monitoring (social media, blogs, forums - with verification)
		\end{itemize}
		
		\textbf{Trigger Events Requiring Unscheduled Refresh:}
		\begin{itemize}
			\item Change in ownership structure (new UBO greater than 25 percent)
			\item Change in business activity (e.g., retail to trade finance)
			\item Geographic risk change (country added to FATF Grey List)
			\item Adverse media hit (name appears in corruption investigation)
			\item PEP designation change (customer becomes PEP)
			\item Unusual transaction pattern (ATL or BTL violation)
			\item Law enforcement inquiry (314(a) request, subpoena, search warrant)
			\item New sanctions designation (customer appears on OFAC SDN list)
		\end{itemize}
		
		\regcite{FinCEN CDD Rule (2016) requires banks to understand nature and purpose of relationships and update on a risk basis.}
		\redflag{Waiting for scheduled 3-year refresh after a material change is a control failure.}
	\end{frame}
	
	\begin{frame}{Prioritization of CDD Refresh Using AI}
		\begin{block}{Risk-Based Prioritization for KYC Refresh}
			AI and machine learning can help prioritize which customers need immediate refresh versus those that can wait for scheduled review.
		\end{block}
		
		\textbf{Factors for AI-Based Prioritization:}
		\begin{itemize}
			\item Time since last refresh (older reviews prioritized higher)
			\item Risk tier (high-risk customers prioritized over low-risk)
			\item Transaction activity volume and velocity (sudden increases trigger review)
			\item Changes in customer behavior patterns (deviation from baseline)
			\item Adverse media hits (severity score and recency)
			\item Sanctions or PEP list changes (new matches or close matches)
			\item Geographic risk changes (customer's country added to monitoring list)
			\item Product usage changes (customer started using high-risk products)
		\end{itemize}
		
		\textbf{Benefits of AI Prioritization:}
		\begin{itemize}
			\item Resources focused on highest-risk customers first
			\item Reduced manual effort for low-risk customers
			\item Faster identification of risk changes
			\item Ability to refresh thousands of customers based on risk scoring
			\item Defensible prioritization for regulators
		\end{itemize}
		
		\redflag{Prioritization methodology must be documented and defensible to regulators.}
	\end{frame}
	
	\begin{frame}{Adverse Media Screening - Technical Implementation}
		\begin{block}{Continuous Adverse Media Monitoring}
			Adverse media tools automatically scrape, classify, and match negative news against customer databases.
		\end{block}
		
		\textbf{How Adverse Media Tools Work (Technical):}
		\begin{enumerate}
			\item Data Ingestion: Web scraping (news, government sites, court records) plus API feeds (Reuters, Bloomberg, regulatory databases)
			\item NLP Processing: Named Entity Recognition extracts persons, organizations, locations. Sentiment analysis classifies negative vs. neutral or positive.
			\item Entity Resolution: Matches extracted names to customer database using fuzzy logic (overcomes misspellings, middle names, initials)
			\item Risk Scoring: Assigns severity (1-10) based on article source, crime type, conviction status, recency
			\item Alert Generation: Creates case for compliance review if score exceeds configured threshold
		\end{enumerate}
		
		\textbf{Credible Sources vs. Non-Credible:}
		\begin{itemize}
			\item Credible: Official gazettes, court records, major news (Reuters, AP, FT, WSJ, BBC), regulatory actions, Interpol notices
			\item Not credible: Unverified social media, anonymous blogs, competitor accusations, unmoderated comment sections
		\end{itemize}
		
		\examfocus{Exam tests distinction between credible adverse media (actionable) and non-credible (not actionable alone).}
	\end{frame}
	
	\begin{frame}{Adverse Media - Actionable Without Conviction Case Study}
		\scenario{CEO Named in Corruption Inquiry - No Conviction}
		An EDD analyst discovers through OSINT that a corporate customer's CEO was named (not charged, not convicted) in a foreign parliamentary corruption inquiry 18 months ago. Separately, the customer's primary counterparty's director appears on an Interpol Red Notice. The compliance officer argues that since no conviction has occurred, the adverse media is inconclusive and no action is required.
		
		\textbf{What the Exam Tests:} Reasonable suspicion, not conviction, is the standard for AML action.
		
		\textbf{How to Analyze (Exam Logic):}
		\begin{enumerate}
			\item OSINT findings from parliamentary inquiry constitute credible adverse media
			\item Conviction is NOT required for AML action (reasonable suspicion standard)
			\item Interpol Red Notice alone justifies immediate SAR filing (international wanted person)
			\item Findings require escalation and potential EDD, relationship reassessment
			\item No conviction is not a defense for ignoring credible adverse media
		\end{enumerate}
		
		\trap{No conviction means the adverse media is inconclusive and can be ignored.}
		\correct{Correct: OSINT findings do not require a conviction to be actionable. Reasonable suspicion is the standard for AML risk assessment.}
	\end{frame}
	
	\begin{frame}{Network Analysis - Centrality Metrics for Money Mule Detection}
		\begin{block}{Graph Theory in AML}
			Network analysis models accounts as nodes and transactions as edges, then calculates centrality metrics to identify suspicious patterns.
		\end{block}
		
		\textbf{Key Centrality Metrics and AML Application:}
		\begin{itemize}
			\item Degree Centrality: Number of direct connections. High degree equals many counterparties. Suspicious for money mules (many incoming small payments).
			\item Betweenness Centrality: Number of shortest paths passing through a node. High betweenness equals critical intermediary. Suspicious for layering accounts.
			\item Eigenvector Centrality: Connected to other high-centrality nodes. Suspicious for core conspiracy members or masterminds.
			\item PageRank (Google algorithm): Importance based on inbound connections. Finds central players in money laundering network.
		\end{itemize}
		
		\textbf{Clustering Algorithms:}
		\begin{itemize}
			\item Louvain community detection: Identifies groups of accounts with dense internal connections (e.g., money mule rings)
			\item Connected components: Accounts that transact only within a closed group (no external economic rationale)
		\end{itemize}
		
		\examfocus{Exam expects you to know that network analysis detects patterns individual monitoring cannot see.}
	\end{frame}
	
	\begin{frame}{Network Analysis Example - Money Mule Ring Detection}
		\scenario{Network Analysis Reveals 23-Account Money Mule Ring}
		A bank's network analysis tool identifies a cluster of 23 accounts with: (1) all opened within a 48-hour window; (2) all received small P2P transfers (\$200-\$500) from over 100 different external senders; (3) all then aggregated funds and wired to the same offshore shell company; (4) shared IP addresses (VPN through secrecy haven); (5) shared mobile phone number for verification. Individual transaction monitoring flagged zero accounts as suspicious because each account's activity appeared isolated.
		
		\textbf{What the Exam Tests:} Network analysis detects coordinated patterns that individual account monitoring misses.
		
		\textbf{How to Analyze (Exam Logic):}
		\begin{enumerate}
			\item Individual monitoring cannot see cross-account patterns
			\item Network analysis reveals the coordinated activity across 23 accounts
			\item The pattern is classic money mule network (layering, consolidation, offshore exit)
			\item Shared IP addresses and phone numbers are critical indicators (tying accounts together)
			\item Bank should file SARs on all 23 accounts
			\item Freeze remaining funds and exit relationships
		\end{enumerate}
		
		\trap{Individual account monitoring is sufficient; network analysis is optional.}
		\correct{Correct: Network analysis is essential for detecting coordinated criminal networks that individual account monitoring cannot identify. The pattern requires SAR filing.}
	\end{frame}
	
	\begin{frame}{Payment Screening - SWIFT MT Fields Required for Screening}
		\begin{block}{SWIFT MT Message Types}
			Most cross-border wires use SWIFT MT (MT103 customer transfer, MT202 interbank transfer).
		\end{block}
		
		\textbf{Key Fields and Screening Requirements:}
		\begin{longtable}{p{0.2\textwidth}|p{0.35\textwidth}|p{0.4\textwidth}}
			\hline
			\rowcolor{lightgray}
			\textbf{Field} & \textbf{Content} & \textbf{Screening Requirement} \\
			\hline
			Field 50 (Ordering Customer) & Name and address of originator & Full screening against sanctions, PEP, adverse media \\
			\hline
			Field 52 (Ordering Institution) & Originating bank & Sanctions screening of bank \\
			\hline
			Field 56 (Intermediary Institution) & Correspondent bank & Sanctions screening \\
			\hline
			Field 57 (Account With Institution) & Beneficiary bank & Sanctions screening \\
			\hline
			Field 59 (Beneficiary Customer) & Name and address of beneficiary & Full screening against sanctions, PEP, adverse media \\
			\hline
			Field 70 (Remittance Information) & Free text payment purpose & NLP screening for suspicious keywords (e.g., discreet, cash, shell) \\
			\hline
			Field 72 (Sender to Receiver Info) & Free text instructions & Screening for evasion patterns \\
			\hline
		\end{longtable}
		
		\redflag{Free text fields (Field 70, 72) are often used to evade screening with misspellings or coded language.}
		\techterm{ISO 20022 (successor to SWIFT MT) provides richer structured data, reducing free text evasion.}
	\end{frame}
	
	\begin{frame}{Blockchain Transaction Screening - Technical Deep Dive}
		\begin{block}{How Blockchain Analytics Work}
			Blockchain surveillance tools map transactions on public ledgers (Bitcoin, Ethereum, etc.) to identify illicit activity.
		\end{block}
		
		\textbf{Technical Capabilities:}
		\begin{itemize}
			\item Address Clustering: Groups addresses owned by the same entity using heuristic analysis (common input ownership, change address patterns, peeling chains)
			\item Transaction Graph Analysis: Follows funds through multiple hops (5-10 plus hops) to trace origin or destination
			\item Risk Scoring: Assigns risk score (0-100) to addresses based on exposure to known illicit sources (darknet markets, mixers, ransomware wallets, sanctioned addresses)
			\item Real-time Screening: Alerts when customer sends or receives to or from high-risk address
			\item Travel Rule Compliance: Originator and beneficiary information attached to transactions (required for regulated VASPs)
		\end{itemize}
		
		\textbf{Major Blockchain Analytics Vendors:}
		\begin{itemize}
			\item Chainalysis: KYT (Know Your Transaction) for real-time screening, Reactor for investigation
			\item Elliptic: Navigator for visual investigation, Lens for real-time risk scoring
			\item TRM Labs: Transaction monitoring, sanctions screening, case management
			\item CipherTrace: (now Mastercard) Cross-border crypto intelligence
		\end{itemize}
		
		\redflag{Privacy coins (Monero, Zcash) cannot be fully traced - many exchanges delist them entirely.}
	\end{frame}
	
	\begin{frame}{FATF Travel Rule - Detailed VASP Obligations}
		\begin{block}{FATF Recommendation 16 (Updated 2019)}
			Virtual Asset Service Providers (VASPs) must share originator and beneficiary information for transfers above threshold (typically \$1,000 USD or EUR).
		\end{block}
		
		\textbf{Required Information for Each Transfer:}
		\begin{itemize}
			\item Originator's full name
			\item Originator's account number (wallet address or account reference)
			\item Originator's address, national ID number, or date of birth (at least one)
			\item Beneficiary's full name
			\item Beneficiary's account number (wallet address)
		\end{itemize}
		
		\textbf{Originating VASP Obligations:}
		\begin{itemize}
			\item Collect and verify required information before transfer
			\item Transmit information securely to receiving VASP
			\item Maintain records for at least 5 years
			\item Immediately freeze and report suspicious transfers
		\end{itemize}
		
		\textbf{Receiving VASP Obligations:}
		\begin{itemize}
			\item Obtain required information from sending VASP
			\item Take action if information not provided (request, delay, block, enhanced scrutiny)
			\item Cannot accept transfers without required information (absolute obligation)
			\item File SAR if information cannot be obtained and transfer is suspicious
		\end{itemize}
		
		\redflag{Receiving VASP cannot simply ignore missing originator data. Not my problem is a violation.}
	\end{frame}
	
	\begin{frame}{Travel Rule Compliance Failure - Exam Scenario}
		\scenario{VASP Receives Transfer with No Originator Information}
		A registered VASP receives a transfer of 4.2 Bitcoin (approximately \$180,000) from another VASP with no originator or beneficiary information attached. The receiving VASP's compliance officer argues that because the customer has already passed KYC, the missing originator data is the sending VASP's problem. The customer states the funds are a personal Bitcoin investment.
		
		\textbf{What the Exam Tests:} Receiving VASP obligations under the Travel Rule.
		
		\textbf{How to Analyze (Exam Logic):}
		\begin{enumerate}
			\item Receiving VASP has independent obligation to obtain required information (FATF Recommendation 16)
			\item Missing information requires follow-up action (request from sending VASP, delay transaction, enhanced scrutiny)
			\item \$180,000 transfer is well above typical \$1,000 threshold
			\item Customer's KYC at receiving VASP does NOT substitute for originator data (different person)
			\item If sending VASP cannot provide information, consider blocking and SAR filing
		\end{enumerate}
		
		\trap{The Travel Rule imposes obligations only on the originating VASP.}
		\correct{Correct: The receiving VASP has an obligation to obtain originator and beneficiary information. If not provided, the receiving VASP must take follow-up action, including enhanced scrutiny and potential SAR filing.}
	\end{frame}
	
	% ============================================================
	% SECTION 5: AI AND MACHINE LEARNING
	% ============================================================
	\section{AI and Machine Learning}
	
	\begin{frame}{Machine Learning Types for AML - Detailed Taxonomy}
		\begin{block}{Supervised Learning (Most Common in AML)}
			Model learns from labeled historical data (suspicious vs. legitimate transactions).
		\end{block}
		
		\textbf{Examples:}
		\begin{itemize}
			\item Random Forest, XGBoost, Gradient Boosting (tree-based models - good explainability)
			\item Logistic Regression (simple, very explainable, good baseline)
			\item Neural Networks / Deep Learning (powerful but black box - requires XAI)
		\end{itemize}
		
		\textbf{Pros:} High accuracy if labeled data available, can optimize for recall (catching bad guys)
		\textbf{Cons:} Requires large labeled dataset (SARs), may not detect new typologies
		
		\begin{block}{Unsupervised Learning (Anomaly Detection)}
			Model finds patterns without labels - flags outliers as potential anomalies.
		\end{block}
		
		\textbf{Examples:}
		\begin{itemize}
			\item Autoencoders (neural networks that reconstruct input, high reconstruction error equals anomaly)
			\item Isolation Forests (random forest that isolates anomalies)
			\item One-Class SVM (learns boundary around normal data)
			\item K-means clustering (distant points from centroids equals anomalies)
		\end{itemize}
		
		\textbf{Pros:} Can detect novel typologies, no labeled data needed
		\textbf{Cons:} Very high false positives (often 90 percent plus), hard to validate, hard to explain
		
		\examfocus{Hybrid approach: supervised for known typologies, unsupervised for novel anomalies.}
	\end{frame}
	
	\begin{frame}{Model Risk Management (MRM) Framework - SR 11-7}
		\begin{block}{Regulatory Requirements for Model Validation}
			All AML models (transaction monitoring, sanctions screening, fraud detection) are subject to model risk management requirements under OCC 2011-12 and SR 11-7.
		\end{block}
		
		\textbf{Key MRM Components:}
		\begin{enumerate}
			\item Model Inventory: All models documented, owners assigned, risk tier (high, medium, low)
			\item Model Development Documentation: Theoretical basis, assumptions, limitations, data sources
			\item Independent Validation: Before deployment and periodically after (annually for high-risk models)
			\item Performance Monitoring: Ongoing tracking of precision, recall, false positive rate, drift
			\item Model Governance: Policies, procedures, approval authorities, escalation path
		\end{enumerate}
		
		\textbf{Validator Independence Requirements:}
		\begin{itemize}
			\item Validator cannot be involved in model development
			\item Validator must report functionally to independent risk management (not to model development)
			\item External validation for high-risk models (e.g., sanctions screening)
		\end{itemize}
		
		\textbf{Validation Report Must Address:}
		\begin{itemize}
			\item Conceptual soundness (does model logic make sense for AML?)
			\item Data quality (is input data complete and accurate?)
			\item Outcomes analysis (precision, recall, ATL, BTL)
			\item Benchmarking (comparison to alternative models)
			\item Limitations and compensating controls
		\end{itemize}
		
		\redflag{A model that cannot be validated (black box AI) is a regulatory violation.}
	\end{frame}
	
	\begin{frame}{Model Drift - Concept Drift vs. Data Drift}
		\begin{block}{Why ML Models Degrade Over Time}
			Models are trained on historical data. When the real world changes, model performance degrades. This is called model drift.
		\end{block}
		
		\textbf{Concept Drift (Most Dangerous):}
		\begin{itemize}
			\item The relationship between input features and target variable changes
			\item Example: Criminal typology shifts from cash structuring to cryptocurrency layering. Old model learned patterns of cash behavior, no longer relevant.
			\item Detection: Monitor recall rate (capture of new SARs). Decline indicates concept drift.
			\item Remediation: Retrain model with new data including new typology.
		\end{itemize}
		
		\textbf{Data Drift (Feature Distribution Changes):}
		\begin{itemize}
			\item Input data distribution changes, but relationship remains same
			\item Example: Customer transaction volume increases 180 percent after product launch. Model trained on lower volumes now sees many false positives.
			\item Detection: Monitor feature distributions (KS test, Population Stability Index)
			\item Remediation: Recalibrate thresholds or retrain with new data
		\end{itemize}
		
		\textbf{Regulatory Expectation:}
		\begin{itemize}
			\item Annual model validation (minimum)
			\item Trigger-based validation after material changes to products, customer base, or typologies
		\end{itemize}
		
		\examfocus{Exam scenario: model not recalibrated for 3.5 years despite new products and typologies - material deficiency.}
	\end{frame}
	
	\begin{frame}{Model Explainability (XAI) - SHAP and LIME}
		\begin{block}{Why Explainability Matters}
			Regulators require that analysts understand why a transaction was flagged. SAR narratives must include rationale. Black box AI is unacceptable.
		\end{block}
		
		\textbf{SHAP (SHapley Additive exPlanations):}
		\begin{itemize}
			\item Game-theoretic approach: assigns each feature an importance value for the prediction
			\item Output example: Transaction amount \$25,000 contributed +0.8 to risk score; counterparty high-risk jurisdiction contributed +0.5; total 1.3 exceeds threshold 1.0
			\item Analysts can see which features drove the alert
			\item Consistent and locally accurate
		\end{itemize}
		
		\textbf{LIME (Local Interpretable Model-agnostic Explanations):}
		\begin{itemize}
			\item Creates a simple, interpretable model (e.g., linear regression) locally around the prediction
			\item Explains individual predictions, not global model behavior
			\item Works with any black box model
			\item Faster than SHAP but less stable
		\end{itemize}
		
		\textbf{Counterfactual Explanations:}
		\begin{itemize}
			\item If the transaction amount had been \$19,000 instead of \$25,000, it would not have been flagged
			\item Useful for customer disputes and internal understanding
		\end{itemize}
		
		\redflag{Black box AI with no explainability will be rejected by regulators regardless of accuracy.}
	\end{frame}
	
	\begin{frame}{Unsupervised Anomaly Detection - Autoencoder Example}
		\begin{block}{How Autoencoders Detect Anomalies}
			Neural network trained to reconstruct normal transactions. High reconstruction error equals anomaly (potential suspicious activity).
		\end{block}
		
		\textbf{Technical Architecture:}
		\begin{itemize}
			\item Encoder: Compresses input (e.g., 100 transaction features) into lower-dimensional representation (e.g., 10 latent variables)
			\item Decoder: Reconstructs original input from latent representation
			\item Training: Minimize reconstruction error on normal transactions only
		\end{itemize}
		
		\textbf{Inference (Monitoring):}
		\begin{itemize}
			\item Feed transaction through autoencoder
			\item Calculate Mean Squared Error (MSE) between input and reconstruction
			\item If MSE greater than threshold, flag as anomaly
		\end{itemize}
		
		\textbf{Pros and Cons:}
		\begin{itemize}
			\item Pro: No labeled data required, can detect novel typologies
			\item Con: Very high false positives (90 percent plus)
			\item Con: Hard to explain why transaction was flagged
			\item Con: Requires threshold tuning (challenging without labels)
		\end{itemize}
		
		\examfocus{Exam tests that unsupervised models still require independent validation, despite being supplemental.}
	\end{frame}
	
	\begin{frame}{Precision, Recall, F1 Score, AUC-ROC for AML}
		\begin{block}{Key Performance Metrics for AML Models}
			Regulators expect banks to track and report these metrics regularly.
		\end{block}
		
		\begin{longtable}{p{0.25\textwidth}|p{0.35\textwidth}|p{0.35\textwidth}}
			\hline
			\rowcolor{lightgray}
			\textbf{Metric} & \textbf{Formula} & \textbf{AML Interpretation} \\
			\hline
			Precision & TP / (TP + FP) & Of all alerts, how many were true positives? Low precision equals many false positives. \\
			\hline
			Recall (Sensitivity) & TP / (TP + FN) & Of all suspicious activity, how many generated alerts? Low recall equals missed suspicious activity. \\
			\hline
			F1 Score & 2 * (Precision * Recall) / (Precision + Recall) & Harmonic mean - balances precision and recall. \\
			\hline
			Specificity & TN / (TN + FP) & Of all legitimate activity, how many did NOT generate alerts? Low specificity equals many false positives. \\
			\hline
			AUC-ROC & Area Under Receiver Operating Characteristic Curve & Model's ability to distinguish suspicious from legitimate (0.5 equals random, 1.0 equals perfect). \\
			\hline
		\end{longtable}
		
		\textbf{Target Ranges for AML (Typical Industry Benchmarks):}
		\begin{itemize}
			\item Recall: Greater than 80 percent (regulators prioritize detection)
			\item Precision: 10-30 percent (many false positives accepted if recall is high)
			\item F1 Score: 0.2-0.4 (typical for AML due to class imbalance)
			\item AUC-ROC: Greater than 0.85 (good discrimination)
		\end{itemize}
		
		\redflag{Regulators are less concerned with low precision than low recall. Do not sacrifice recall for precision.}
	\end{frame}
	
	\begin{frame}{AI in Investigations - Automated Narrative Generation}
		\begin{block}{AI Tools to Support AML Investigators}
			Beyond alert generation, AI and ML can significantly enhance the investigation workflow.
		\end{block}
		
		\textbf{AI-Powered Investigation Tools:}
		\begin{itemize}
			\item Automated Data Aggregation: Pulls KYC, transactions, sanctions hits, adverse media into single case file (minutes versus hours)
			\item Entity Resolution: Connects related accounts, addresses, phone numbers across systems
			\item Risk Scoring of Alerts: Prioritizes highest-risk alerts for investigator review
			\item Automated Narrative Generation: Drafts SAR narratives from structured case data (analyst reviews and edits - not auto-filed)
			\item Link Analysis Visualization: Automatically maps relationships between entities
			\item Predictive Case Closure: Suggests likely false positive based on historical patterns (with human validation)
		\end{itemize}
		
		\textbf{Key Exam Distinction:}
		\begin{itemize}
			\item AI can AUTOMATE data collection and prioritization
			\item AI should NOT replace investigator judgment for SAR decisions
			\item Investigators must always review and validate AI suggestions
			\item A SAR filed solely on AI recommendation without human analysis is a regulatory risk
		\end{itemize}
		
		\redflag{A SAR filed solely on an AI recommendation without human analysis is unacceptable.}
	\end{frame}
	
	% ============================================================
	% SECTION 6: INVESTIGATION TOOLS AND OSINT
	% ============================================================
	\section{Investigation Tools and OSINT}
	
	\begin{frame}{Case Management System Requirements}
		\begin{block}{Core Functionality for AML Investigators}
			Case management systems (CMS) support the investigation workflow from alert to SAR filing.
		\end{block}
		
		\textbf{Must-Have CMS Capabilities:}
		\begin{itemize}
			\item Alert Intake: Auto-creation of cases from monitoring system alerts, with priority scoring (high, medium, low)
			\item Data Aggregation: Automatically pull KYC, transactions, sanctions hits, adverse media, network links into single case file
			\item Link Analysis Visualization: Graph display of entity relationships, transaction flows
			\item Timeline Analysis: Chronological view of customer activity (list view plus graph)
			\item Document Management: Store investigation notes, evidence, screenshots, emails
			\item SAR Generation: Auto-populate SAR or STR forms from case data (populated by investigator, not auto-filed)
			\item Audit Trail: Log all actions (who viewed what, when, what decisions made) - immutable
			\item Collaboration Tools: Notes, tasks, assignments, supervisor approval workflow
			\item Reporting Dashboards: Management view of queue size, aging, SAR volume, analyst performance
		\end{itemize}
		
		\redflag{Audit trail is mandatory - regulators will request logs of who accessed which cases.}
	\end{frame}
	
	\begin{frame}{Open Source Intelligence (OSINT) - Credible Sources Detailed}
		\begin{block}{Credible OSINT Sources (Actionable for AML)}
			Only information from authoritative, verifiable sources should be used in AML investigations.
		\end{block}
		
		\textbf{Credible Sources by Category:}
		\begin{itemize}
			\item Government Official Sources: Sanctions lists (OFAC, UN, EU, UK HMT), corporate registries (Companies House, SEC EDGAR), court records (PACER), government gazettes
			\item Law Enforcement: Interpol Red Notices (www.interpol.int), FBI Most Wanted, Europol wanted lists
			\item Regulatory: FinCEN enforcement actions, SEC enforcement, FCA fines, OCC consent orders
			\item Investigative Journalism: ICIJ (Offshore Leaks, Panama Papers, Pandora Papers), OCCRP, Bellingcat, Reuters Investigates
			\item Major News: Reuters, Bloomberg, Financial Times, Associated Press, Wall Street Journal, BBC
			\item Commercial Databases: Dun and Bradstreet, LexisNexis, Dow Jones Factiva, Sayari
		\end{itemize}
		
		\textbf{Non-Credible Sources (Not Actionable Alone):}
		\begin{itemize}
			\item Unverified social media (Twitter, Facebook, LinkedIn without verification)
			\item Anonymous blogs (Medium, WordPress with no author attribution)
			\item Unmoderated forums (Reddit, 4chan)
			\item Competitor allegations without evidence
			\item Personal websites with no attribution
		\end{itemize}
		
		\examfocus{Exam tests distinction between credible and non-credible OSINT sources.}
	\end{frame}
	
	% ============================================================
	% SECTION 7: PRIVACY AND DATA PROTECTION
	% ============================================================
	\section{Privacy and Data Protection}
	
	\begin{frame}{Privacy-Enhancing Technologies (PETs) for AML}
		\begin{block}{Balancing AML Obligations and Privacy Rights}
			Technology can help comply with both AML and privacy laws (GDPR, CCPA, etc.).
		\end{block}
		
		\textbf{PETs Applicable to AML:}
		\begin{itemize}
			\item Data Minimization: Collect only data required for AML (not extraneous personal data). GDPR Article 5(1)(c).
			\item Anonymization: Remove all personal identifiers where not needed (e.g., aggregated reports). Cannot be reversed.
			\item Pseudonymization: Replace identifiers with tokens (e.g., customer ID replaced with token). Reversible with key. GDPR encourages pseudonymization.
			\item Encryption: Protect data at rest (AES-256) and in transit (TLS 1.3). Required for data protection laws.
			\item Access Controls: Role-based access (RBAC) - investigators see only data needed for their cases. Log all access.
			\item Audit Logs: Track all access to sensitive AML data. Required for breach notification.
			\item Secure Data Sharing: Encrypted channels for 314(b) sharing, FIU reporting (GoAML, secure FTP).
		\end{itemize}
		
		\redflag{AML obligations override privacy rights for required data retention (e.g., 5-year recordkeeping).}
	\end{frame}
	
	\begin{frame}{GDPR and AML - Right to Erasure Exception}
		\scenario{Bank Receives Right to Erasure Request for SAR Customer}
		A European bank receives a GDPR right to erasure request from a customer demanding deletion of all personal data, including KYC documents and transaction records. The customer was the subject of a SAR filed 18 months ago. The Data Privacy Officer argues GDPR requires deletion. The AML Officer argues AML recordkeeping requires retention (minimum 5 years).
		
		\textbf{What the Exam Tests:} GDPR exception for legal obligations (AML retention).
		
		\textbf{How to Analyze (Exam Logic):}
		\begin{enumerate}
			\item GDPR Article 17(3)(b) exempts data retention required by Union or Member State law
			\item AML directives (4AMLD, 5AMLD) require 5-year retention of records
			\item Bank can lawfully refuse to delete AML-required records
			\item SAR confidentiality (non-tipping off) also prohibits acknowledging SAR existence to customer
		\end{enumerate}
		
		\trap{GDPR's right to erasure is absolute and overrides all other laws.}
		\correct{Correct: AML retention requirements are a legal exception to GDPR erasure rights. The bank may refuse deletion for records it is required to retain.}
	\end{frame}
	
	% ============================================================
	% SECTION 8: TOOL SELECTION AND INTEGRATION
	% ============================================================
	\section{Tool Selection and Integration}
	
	\begin{frame}{Tool Selection Criteria - Comprehensive Framework}
		\begin{block}{Evaluating AML Technology Vendors}
			Selecting the right tool requires balancing multiple factors; cheapest is rarely best.
		\end{block}
		
		\textbf{Key Selection Criteria with Sub-Questions:}
		\begin{itemize}
			\item Risk Alignment: Does the tool address your specific ML or TF risks (trade finance, crypto, correspondent banking, etc.)? Request vendor typology coverage list.
			\item Regulatory Acceptance: Has the vendor been examined by regulators (OCC, FinCEN, FCA)? Do peer institutions use it? Any public enforcement actions against vendor?
			\item Integration Complexity: APIs available (REST, GraphQL)? Data format compatibility with core banking system? Proof of concept with your data required.
			\item Performance Metrics: Request vendor's false positive rate, recall rate, precision on independent test data (not vendor's own marketing numbers).
			\item Explainability: Can the tool explain WHY it flagged a transaction? Required for AI models (SHAP, LIME, rule extraction).
			\item Total Cost of Ownership (TCO): License fees plus implementation plus integration plus training plus ongoing maintenance plus hardware or cloud costs.
			\item Vendor Stability: Financial health, support responsiveness, update frequency, data privacy compliance (GDPR, CCPA).
			\item Scalability: Can it handle 2x or 5x current transaction volume without performance degradation? Stress test results?
		\end{itemize}
		
		\redflag{Failure to conduct proper due diligence on a vendor is a regulatory deficiency.}
	\end{frame}
	
	\begin{frame}{Integration Architecture - Connecting AFC Tools}
		\begin{block}{Integration Risk Points (Exam Scenarios)}
			A tool is only as good as its data feeds. Poor integration equals poor detection.
		\end{block}
		
		\textbf{Systems That Must Integrate:}
		\begin{itemize}
			\item Core Banking System: Customer master, account balances, transaction history
			\item Payment Systems: SWIFT, ACH, Fedwire, SEPA, crypto wallets (via API)
			\item CRM: Customer contact info, relationship manager assignments
			\item Case Management: Workflow integration (alert to case to SAR)
			\item Data Lake or Warehouse: Historical data for model training and validation
			\item Sanctions Lists: Real-time API feeds for list updates (avoid manual downloads)
		\end{itemize}
		
		\textbf{Integration Risk Points (Common Exam Scenarios):}
		\begin{itemize}
			\item Data format mismatches: CSV vs. JSON vs. XML vs. proprietary formats
			\item Data latency: Transaction at 9:00 AM, monitoring system sees it at 9:00 PM (criminal window)
			\item Silent feed failures: No error message, data just missing (most dangerous)
			\item API rate limiting: Vendor caps queries at 1,000 per second, but bank has 5,000 per second equals queued or dropped transactions
			\item Version compatibility: Core system upgrade breaks integration for weeks
		\end{itemize}
		
		\textbf{Best Practices:}
		\begin{itemize}
			\item Always require a proof of concept (POC) with your actual data
			\item Test data feeds thoroughly before go-live (including failure scenarios)
			\item Establish reconciliation between source and target
			\item Document integration architecture and data flows
		\end{itemize}
		
		\redflag{API rate limiting is a common hidden risk - test at peak volume.}
	\end{frame}
	
	% ============================================================
	% SECTION 9: OPERATIONAL EFFECTIVENESS AND EFFICIENCY
	% ============================================================
	\section{Operational Effectiveness and Efficiency}
	
	\begin{frame}{Operational Effectiveness vs. Efficiency - Regulatory View}
		\begin{block}{Balancing Two Competing Priorities}
			Regulators prioritize effectiveness (catching bad activity) over efficiency (cost per alert).
		\end{block}
		
		\textbf{Measuring Operational Effectiveness (Regulator Focus):}
		\begin{itemize}
			\item False negative rate (missed suspicious activity) - should be less than 5 percent
			\item SAR quality (actionable intelligence for law enforcement) - percentage of SARs that result in investigation or prosecution
			\item Investigation timeliness (alerts reviewed within target timeframes) - e.g., 48 hours for high-priority alerts
			\item Alert-to-SAR conversion rate (meaningful alerts versus noise) - 2-5 percent is typical
			\item Audit findings - fewer findings indicate better controls
			\item Back-testing results - historical SAR capture rate
		\end{itemize}
		
		\textbf{Measuring Operational Efficiency (Bank Management Focus):}
		\begin{itemize}
			\item Cost per alert reviewed (typical \$25-\$75)
			\item Time per investigation (hours per case)
			\item False positive rate (higher equals less efficient)
			\item Analyst utilization rate (percentage of time on productive investigation)
			\item SAR filing rate relative to peer institutions (too low or too high both scrutinized)
		\end{itemize}
		
		\redflag{Efficiency improvements that increase false negatives (e.g., raising thresholds arbitrarily) are unacceptable.}
	\end{frame}
	
	\begin{frame}{Risk-Based Resource Prioritization - Investment Mapping}
		\begin{block}{Aligning Investment with Risk Assessment}
			The risk-based approach means investing most heavily where risk is highest.
		\end{block}
		
		\textbf{High-Risk Areas Requiring More Investment:}
		\begin{itemize}
			\item Correspondent Banking (highest risk): Enhanced due diligence (Wolfsberg CBDDQ), on-site inspections, dedicated monitoring team, automated screening of nested accounts
			\item Private Banking and Wealth Management (PEP risk): Enhanced CDD, source of wealth verification, senior management approval, relationship manager oversight
			\item Trade Finance (TBML vulnerability): Automated invoice matching, country risk scoring, counterparty screening, dedicated TBML analysts
			\item Virtual Assets (emerging typologies): Blockchain analytics (Chainalysis or Elliptic), Travel Rule compliance, mixer and darknet monitoring
			\item Cross-Border Payments (jurisdictional risk): Enhanced payment screening, jurisdiction risk scoring, correspondent bank restrictions
			\item Cash-Intensive Businesses (structuring risk): Aggregated cash reporting, geographic profiling, branch-level monitoring
		\end{itemize}
		
		\textbf{Lower-Risk Areas for Lighter Investment (Simplified Due Diligence):}
		\begin{itemize}
			\item Domestic retail payroll accounts (low transaction velocity, predictable patterns)
			\item Low-balance consumer accounts (less than \$5,000 balance)
			\item Government benefit disbursement accounts (limited usage)
			\item Regulated public companies (subject to disclosure requirements)
		\end{itemize}
		
		\redflag{Applying the same resources to all customers violates the risk-based approach.}
	\end{frame}
	
	% ============================================================
	% SECTION 10: FINAL REVIEW - TOP CONCEPTS
	% ============================================================
	\section{Final Review}
	
	\begin{frame}{Domain D Summary - Top 50 Most Tested Concepts (Part 1)}
		\begin{multicols}{2}
			\begin{enumerate}
				\item Silent data failure equals monitoring gap
				\item Reconciliation controls prevent undetected data loss
				\item Fuzzy logic required for sanctions screening
				\item False negatives worse than false positives
				\item Model drift requires regular recalibration
				\item Explainability required for AI models (XAI)
				\item E-KYC requires liveness detection
				\item Network analysis detects money mule rings
				\item Travel Rule applies to both sending and receiving VASPs
				\item Adverse media does NOT require conviction
				\item Data quality is compliance responsibility
				\item Perpetual KYC is superior to periodic refresh
				\item Trigger-based refresh required for material changes
				\item ATL and BTL testing validates scenario effectiveness
				\item Rules-based plus ML hybrid is best practice
				\item Unsupervised models still require validation
				\item SWIFT fields requiring screening: 50, 52, 56, 57, 59, 70
				\item Blockchain analytics required for VASPs
				\item Privacy coins and mixers require EDD
				\item Receiving VASP cannot ignore missing Travel Rule data
				\item GDPR has AML retention exception (Article 17(3)(b))
				\item OSINT must use credible sources only
				\item Tool selection must consider integration and scalability
				\item Effectiveness prioritized over efficiency
				\item Risk-based approach drives resource allocation
			\end{enumerate}
		\end{multicols}
	\end{frame}
	
	\begin{frame}{Domain D Summary - Top 50 Most Tested Concepts (Part 2)}
		\begin{multicols}{2}
			\begin{enumerate}
				\setcounter{enumi}{25}
				\item Model risk management requires independent validation
				\item Sanctions screening requires fuzzy logic, not exact matching
				\item Over-filtering increases false negatives
				\item Scenario coverage must address all typologies
				\item Investigation tools require audit trails
				\item Biometrics alone insufficient without liveness
				\item MRZ validation for identity documents
				\item Hash totals for data reconciliation
				\item Concept drift versus data drift
				\item SHAP and LIME for model explainability
				\item Precision, recall, F1, AUC-ROC metrics
				\item Autoencoders for anomaly detection
				\item High betweenness centrality equals layering account
				\item FATF digital ID assurance levels (low, medium, high)
				\item Silent feed failures are most dangerous
				\item NULL Customer ID equals unlinked transactions
				\item API rate limiting causes missed transactions
				\item Proof of concept required before vendor contract
				\item Analyst fatigue from high false positives
				\item Board reporting on data quality metrics
				\item Non-tipping off for SAR customers
				\item Interpol Red Notice equals immediate SAR consideration
				\item Correspondent banking highest risk
				\item Nested account activity requires EDD
				\item Shell bank prohibition (USA PATRIOT Act)
			\end{enumerate}
		\end{multicols}
	\end{frame}
	
	\begin{frame}{Critical Red Flags - Domain D Quick Reference}
		\begin{multicols}{2}
			\begin{itemize}
				\item 34 percent missing data for 14 months -> Data failure
				\item NULL Customer ID on transactions -> Unlinked transactions
				\item No reconciliation controls -> Undetected gaps
				\item 77 percent false negative rate -> Material deficiency
				\item No recalibration for 3.5 years -> Model drift
				\item Black box AI with no explainability -> Governance failure
				\item Exact-string matching only -> Sanctions screening failure
				\item No liveness detection -> Spoofing vulnerability
				\item No adverse media review -> Missed risk indicators
				\item No network analysis -> Missed mule rings
				\item Missing Travel Rule data accepted -> VASP violation
				\item No ATL or BTL testing -> Unvalidated scenarios
				\item Only annual refresh after material changes -> Control gap
				\item Same onboarding for all customers -> RBA violation
				\item No independent model validation -> MRM failure
				\item API rate limiting untested -> Transaction loss
				\item No hash reconciliation -> Data corruption undetected
				\item 24+ hour data latency -> Criminal window
			\end{itemize}
		\end{multicols}
	\end{frame}
	
	\begin{frame}{Exam Day Reminders for Domain D}
		\begin{itemize}
			\item Domain D is 20 percent of CAMS exam - do not underestimate
			\item Data quality is a COMPLIANCE responsibility, not just IT
			\item Silent data failures (no error messages) are the most dangerous
			\item FALSE NEGATIVES are worse than false positives (regulator priority)
			\item Fuzzy logic is REQUIRED for sanctions screening (exact matching insufficient)
			\item ML models require explainability (SHAP or LIME) and ongoing validation
			\item Model drift requires regular recalibration (annual minimum)
			\item Travel Rule applies to BOTH sending and receiving VASPs
			\item E-KYC without liveness detection can be spoofed (deepfakes)
			\item Adverse media findings do NOT require conviction (reasonable suspicion)
			\item Network analysis detects patterns individual monitoring misses
			\item GDPR has an AML retention exception (Article 17(3)(b))
			\item Credible OSINT sources: government, courts, reputable journalism
			\item Privacy coins and mixers require EDD (high risk)
			\item Effectiveness is greater than efficiency (regulator priority)
			\item Risk-based approach drives resource allocation
			\item Always require a proof of concept (POC) before signing vendor contract
			\item Test reconciliation controls and API rate limits at peak volume
		\end{itemize}
		\vspace{0.3cm}
		\centering
		\greennote{Good luck on the exam!}
	\end{frame}
	
	
	
	% ============================================================
	% SECTION: ONGOING CONTROLS - TRANSACTION MONITORING
	% ============================================================
	
	\begin{frame}{Ongoing Controls: Transaction Monitoring Overview}
		\begin{block}{Core Definition}
			Transaction monitoring is the primary detective control for identifying suspicious activity after customer onboarding. It continuously screens transactions for patterns indicative of money laundering or terrorist financing.
		\end{block}
		
		\textbf{Key Components of Transaction Monitoring:}
		\begin{itemize}
			\item Data ingestion from core banking and payment systems
			\item Rules-based scenarios and machine learning models
			\item Customer segmentation for risk-appropriate thresholds
			\item Alert generation and prioritization
			\item Case management and investigation workflow
			\item SAR/STR filing and documentation
		\end{itemize}
		
		\examfocus{Transaction monitoring is the largest operational investment for most banks and the most frequently cited deficiency in regulatory enforcement actions.}
	\end{frame}
	
	\begin{frame}{Traditional Transaction Monitoring Architecture}
		\begin{block}{Standard Data Flow}
			Core Systems -> Data Aggregation -> Scenario Execution -> Alert Generation -> Case Management -> Investigation -> SAR Filing
		\end{block}
		
		\textbf{System Components in Detail:}
		\begin{itemize}
			\item Data Ingestion Layer: Pulls from core banking, payments, CRM, trade finance. Handles batch (nightly) and real-time (Kafka, RabbitMQ) feeds.
			\item Data Transformation Layer: Normalizes formats, maps fields, enriches with customer risk rating and segmentation data.
			\item Scenario Engine: Executes pre-defined rules (SQL, Drools) or ML models. Must scale to millions of transactions daily.
			\item Alert Generation: Creates case in case management system with priority scoring (high, medium, low).
			\item Case Management: Workflow for investigation, evidence collection, narrative writing, approval.
			\item Reporting Engine: Files SAR/STR to FIU, generates management dashboards.
		\end{itemize}
		
		\redflag{Data latency at any point creates a window for criminals to move funds before detection.}
	\end{frame}
	
	\begin{frame}{Traditional Rules-Based Monitoring - How It Works}
		\begin{block}{Deterministic Logic}
			Rules-based monitoring uses pre-defined, deterministic logic written by analysts to flag specific transaction patterns.
		\end{block}
		
		\textbf{Example Rule Structures:}
		\begin{itemize}
			\item IF cash deposits in rolling 30 days > \$9,500 AND < \$10,000 AND number of deposits > 5 THEN generate alert
			\item IF wire transfer to high-risk jurisdiction AND amount > \$25,000 AND no prior history THEN generate alert
			\item IF rapid movement (funds in and out within 24 hours) AND amount > \$50,000 THEN generate alert
		\end{itemize}
		
		\textbf{Advantages of Rules-Based:}
		\begin{itemize}
			\item Completely explainable and auditable
			\item Low false positive rate when well-tuned (5-10 percent)
			\item Simple to implement with SQL or rules engines
			\item No model drift over time
			\item Low computational requirements
		\end{itemize}
		
		\textbf{Disadvantages of Rules-Based:}
		\begin{itemize}
			\item Misses new or evolving typologies
			\item Requires constant manual rule updates
			\item Criminals can learn and game the rules
			\item Cannot detect complex patterns across multiple dimensions
		\end{itemize}
		
		\examfocus{The exam tests that rules-based systems must be updated regularly as typologies evolve.}
	\end{frame}
	
	\begin{frame}{Customer Segmentation - Risk-Based Monitoring}
		\begin{block}{Why Segmentation is Critical}
			Applying the same monitoring rules to all customers creates excessive false positives for low-risk customers and false negatives for high-risk customers.
		\end{block}
		
		\textbf{Typical Segmentation Dimensions:}
		\begin{itemize}
			\item Customer Risk Tier: Low, Medium, High, PEP
			\item Product Type: Deposit accounts, loans, trade finance, crypto
			\item Geographic Risk: Low, Medium, High, Embargoed
			\item Channel: Branch, online, mobile, ATM
			\item Transaction Type: Cash, wire, ACH, check, crypto
			\item Business Line: Retail, Commercial, Private Banking, Correspondent
		\end{itemize}
		
		\textbf{Segmentation Impact on Thresholds:}
		\begin{longtable}{p{0.3\textwidth}|p{0.3\textwidth}|p{0.3\textwidth}}
			\hline
			\rowcolor{lightgray}
			\textbf{Customer Type} & \textbf{Alert Threshold} & \textbf{Review Priority} \\
			\hline
			Low Risk Retail & Higher threshold & Lower priority \\
			\hline
			Medium Risk Retail & Standard threshold & Standard priority \\
			\hline
			High Risk PEP & Lower threshold & High priority \\
			\hline
			Correspondent Bank & Lowest threshold & Highest priority \\
			\hline
		\end{longtable}
		
		\redflag{Failure to segment customers by risk is a violation of the risk-based approach.}
	\end{frame}
	
	\begin{frame}{Customer Segmentation - Implementation Examples}
		\begin{block}{How Segmentation Changes Monitoring Behavior}
			The same transaction amount triggers different responses based on customer segment.
		\end{block}
		
		\textbf{Example: \$15,000 Wire Transfer to High-Risk Country}
		\begin{itemize}
			\item Low-risk retail customer: No alert (threshold is \$25,000)
			\item Medium-risk retail customer: Low priority alert (threshold is \$10,000)
			\item High-risk customer: High priority alert (threshold is \$5,000)
			\item PEP customer: Immediate block pending review (threshold is \$1,000)
		\end{itemize}
		
		\textbf{Segmentation Data Sources:}
		\begin{itemize}
			\item Customer risk rating from onboarding and periodic review
			\item Transaction history and velocity patterns
			\item Product usage and changes over time
			\item Geographic exposure and counterparty risk
			\item Adverse media and PEP status
		\end{itemize}
		
		\textbf{Segmentation Maintenance:}
		\begin{itemize}
			\item Segments must be reviewed at least annually
			\item Customers can move between segments based on behavior
			\item Automated re-segmentation based on trigger events
		\end{itemize}
		
		\examfocus{Segmentation is not a one-time exercise. Customers change risk profiles over time.}
	\end{frame}
	
	\begin{frame}{Scenario Coverage - Mapping to Typologies}
		\begin{block}{Required Scenario Coverage}
			Every money laundering typology relevant to your institution must have corresponding monitoring scenarios.
		\end{block}
		
		\textbf{Common Typologies and Detection Scenarios:}
		\begin{longtable}{p{0.3\textwidth}|p{0.35\textwidth}|p{0.3\textwidth}}
			\hline
			\rowcolor{lightgray}
			\textbf{Typology} & \textbf{Detection Scenario} & \textbf{Key Indicators} \\
			\hline
			Structuring & Cash aggregation just below CTR threshold & Multiple deposits, multiple branches, multiple tellers \\
			\hline
			Layering & Rapid movement through multiple accounts & Funds in and out within 24 hours, 3+ hops \\
			\hline
			Trade-Based ML & Invoice-payment mismatch & Round-dollar amounts, commodity mismatch \\
			\hline
			Shell Companies & No operating expenses but high volume & No payroll, no rent, no utilities \\
			\hline
			PEP Activity & Large unexplained transfers & No source of funds documentation \\
			\hline
			NPO/Terrorist Financing & Rapid cash out to conflict zone & Donations from unrelated parties \\
			\hline
			Cybercrime & Payment to known ransomware wallet & Darknet market exposure \\
			\hline
			Human Trafficking & Small payments to controller account & Hotel and travel purchases \\
			\hline
		\end{longtable}
		
		\redflag{Missing a scenario for a relevant typology is a regulatory deficiency.}
	\end{frame}
	
	\begin{frame}{Scenario Coverage - Gap Analysis}
		\begin{block}{How to Identify Coverage Gaps}
			Regular gap analysis ensures scenarios remain aligned with your risk assessment.
		\end{block}
		
		\textbf{Gap Analysis Methodology:}
		\begin{enumerate}
			\item Review current risk assessment for identified typologies
			\item Map each typology to existing monitoring scenarios
			\item Identify typologies with no or partial coverage
			\item Prioritize gap remediation based on risk level
			\item Develop and test new scenarios for gaps
		\end{enumerate}
		
		\textbf{Common Coverage Gaps (Exam Scenarios):}
		\begin{itemize}
			\item Trade-based money laundering scenarios missing when institution offers trade finance
			\item Cryptocurrency scenarios missing when institution allows crypto transactions
			\item Correspondent banking nested account scenarios missing
			\item Human trafficking indicators not configured
			\item Rapid layering detection missing (funds in and out same day)
		\end{itemize}
		
		\textbf{Regulatory Expectation:}
		\begin{itemize}
			\item Gap analysis performed at least annually
			\item Documentation of coverage decisions
			\item Remediation plan for identified gaps
			\item Board reporting on significant gaps
		\end{itemize}
		
		\examfocus{The exam tests that scenario coverage must be based on the institution's risk assessment.}
	\end{frame}
	
	\begin{frame}{Threshold Setting - Fundamentals}
		\begin{block}{What Thresholds Control}
			Thresholds determine which transactions generate alerts and which are filtered out.
		\end{block}
		
		\textbf{Types of Thresholds:}
		\begin{itemize}
			\item Monetary Thresholds: Dollar amounts (e.g., alert on transactions > \$10,000)
			\item Velocity Thresholds: Frequency within time period (e.g., 5+ deposits in 30 days)
			\item Count Thresholds: Number of transactions (e.g., 10+ wires to same beneficiary)
			\item Percentage Thresholds: Variance from baseline (e.g., 200% increase from average)
			\item Scoring Thresholds: Risk score from ML model (e.g., score > 0.85)
		\end{itemize}
		
		\textbf{Factors Influencing Threshold Setting:}
		\begin{itemize}
			\item Customer risk tier (lower thresholds for higher risk)
			\item Product type (lower thresholds for high-risk products)
			\item Geographic risk (lower thresholds for high-risk jurisdictions)
			\item Historical false positive rates (adjust to optimize)
			\item Regulatory expectations (specific guidance may exist)
			\item Operational capacity (can you review all alerts?)
		\end{itemize}
		
		\redflag{Thresholds set too high create false negatives (missed suspicious activity).}
	\end{frame}
	
	\begin{frame}{Threshold Setting - Calibration Methodology}
		\begin{block}{Data-Driven Threshold Calibration}
			Thresholds must be set based on data analysis, not arbitrary decisions.
		\end{block}
		
		\textbf{Calibration Process:}
		\begin{enumerate}
			\item Analyze historical transaction data by segment
			\item Identify normal patterns and outliers
			\item Run historical SARs through candidate thresholds
			\item Calculate capture rate (ATL) at each threshold
			\item Calculate false positive rate (BTL) at each threshold
			\item Select threshold balancing detection and operational load
		\end{enumerate}
		
		\textbf{Example Calibration Analysis:}
		\begin{longtable}{p{0.2\textwidth}|p{0.25\textwidth}|p{0.25\textwidth}|p{0.25\textwidth}}
			\hline
			\rowcolor{lightgray}
			\textbf{Threshold} & \textbf{Expected Alerts/ month} & \textbf{Capture Rate (ATL)} & \textbf{False Positive Rate} \\
			\hline
			\$5,000 & 12,000 & 98\% & 35\% \\
			\hline
			\$10,000 & 6,500 & 94\% & 18\% \\
			\hline
			\$15,000 & 3,200 & 87\% & 9\% \\
			\hline
			\$20,000 & 1,800 & 76\% & 5\% \\
			\hline
		\end{longtable}
		
		\textbf{Decision Factors:}
		\begin{itemize}
			\item Regulator priority: higher capture rate (even with more alerts)
			\item Operational capacity: can you review 6,500 alerts per month?
			\item Risk tolerance: is 87% capture rate acceptable?
		\end{itemize}
		
		\examfocus{The exam tests that threshold changes must be data-driven and validated.}
	\end{frame}
	
	\begin{frame}{The Over-Filtering Trap - Detailed Analysis}
		\scenario{Compliance Director Raises All Thresholds by 250 Percent}
		A bank generates 8,000 alerts per month. 96 percent are false positives. The compliance director, under pressure to reduce costs, orders all thresholds raised by 250 percent across all scenarios. The analytics team warns of increased false negatives. The director states: We cannot afford to review 8,000 alerts per month. Any missed activity is acceptable as long as we still file some SARs. Regulators later find criminal structuring below the new thresholds went undetected for 18 months.
		
		\textbf{Why This is a Violation:}
		\begin{enumerate}
			\item 250 percent increase is arbitrary - no data analysis performed
			\item Director prioritized cost over effectiveness (inverse of regulatory priority)
			\item Statement any missed activity is acceptable is willful non-compliance
			\item No consideration of risk-based approach (same increase for all segments)
			\item No ATL/BTL testing to validate new thresholds
		\end{enumerate}
		
		\textbf{Proper Approach:}
		\begin{enumerate}
			\item Analyze false positive root causes (improve scenarios, don't just raise thresholds)
			\item Segment customers (lower thresholds for high-risk, higher for low-risk)
			\item Test threshold changes with ATL/BTL before implementation
			\item Monitor false negative rate after changes
			\item Document all threshold decisions with supporting data
		\end{enumerate}
		
		\trap{The director is correctly managing operational efficiency under the risk-based approach.}
		\correct{Correct: Raising thresholds without data analysis and validation creates unacceptable false negative risk.}
	\end{frame}
	
	\begin{frame}{Statistical Testing - ATL (Above Threshold Limit)}
		\begin{block}{ATL Testing Validates Scenario Effectiveness}
			ATL testing proves that transactions which SHOULD generate alerts actually do.
		\end{block}
		
		\textbf{Detailed Methodology:}
		\begin{enumerate}
			\item Compile historical confirmed suspicious cases (SARs filed, EDD findings)
			\item Ensure cases represent all relevant typologies
			\item Run these cases through current production monitoring scenarios
			\item Calculate capture rate = (cases that generated alert) / (total cases)
			\item Expected capture rate: 90-100 percent for well-tuned scenarios
		\end{enumerate}
		
		\textbf{Low Capture Rate Analysis (Below 90 percent):}
		\begin{itemize}
			\item Thresholds too high (over-filtering) - most common cause
			\item Scenario logic too narrow - missing variations of typology
			\item Data quality issues - missing transactions or fields
			\item Typology has evolved - scenario needs update
			\item Wrong customer segmentation - rules don't apply to this segment
		\end{itemize}
		
		\textbf{Regulatory Expectation:}
		\begin{itemize}
			\item ATL testing at least annually
			\item Material changes trigger interim testing
			\item Results documented and presented to management
			\item Remediation required for out-of-tolerance results
		\end{itemize}
		
		\redflag{A 77 percent capture rate (23 percent false negative) on historical SARs is a material deficiency.}
	\end{frame}
	
	\begin{frame}{ATL Testing - Practical Example}
		\begin{block}{Example: Structuring Scenario ATL Test}
			Testing a scenario designed to detect cash structuring just below CTR threshold.
		\end{block}
		
		\textbf{Test Data Set:}
		\begin{itemize}
			\item 100 historical SARs for cash structuring
			\item Each case involved 5-10 deposits of $8,000-$9,900 within 30 days
			\item Across multiple branches and tellers
		\end{itemize}
		
		\textbf{Run Through Current Scenario:}
		\begin{itemize}
			\item Scenario logic: Aggregate cash deposits < $10,000 within rolling 30 days
			\item Current threshold: 8 or more deposits triggers alert
		\end{itemize}
		
		\textbf{Results:}
		\begin{itemize}
			\item 82 of 100 cases generated alerts (82 percent capture rate)
			\item 18 cases missed (18 percent false negative rate)
		\end{itemize}
		
		\textbf{Root Cause Analysis:}
		\begin{itemize}
			\item Missed cases had only 6-7 deposits (below 8 threshold)
			\item Criminals adapted to 8-deposit threshold
		\end{itemize}
		
		\textbf{Remediation:}
		\begin{itemize}
			\item Lower threshold to 5 deposits
			\item Add secondary indicators (multiple branches, multiple tellers)
			\item Re-test: 97 percent capture rate achieved
		\end{itemize}
		
		\examfocus{ATL testing revealed that criminals adapted to the threshold. Scenarios must evolve.}
	\end{frame}
	
	\begin{frame}{Statistical Testing - BTL (Below Threshold Limit)}
		\begin{block}{BTL Testing Validates Scenario Specificity}
			BTL testing proves that transactions which SHOULD NOT generate alerts do not.
		\end{block}
		
		\textbf{Detailed Methodology:}
		\begin{enumerate}
			\item Identify known legitimate, low-risk transactions
			\item Examples: payroll deposits, utility payments, mortgage payments
			\item Run these transactions through current monitoring scenarios
			\item Calculate false positive rate = (alerts generated) / (total test transactions)
			\item Expected false positive rate: 0-10 percent for well-tuned scenarios
		\end{enumerate}
		
		\textbf{High False Positive Rate Analysis (Above 10 percent):}
		\begin{itemize}
			\item Thresholds too low (under-filtering) - too sensitive
			\item Scenario logic too broad (over-inclusive)
			\item Customer segmentation issues (same rules for different risk tiers)
			\item Data quality issues (corrupted fields causing false matches)
			\item Normal behavior incorrectly classified as suspicious
		\end{itemize}
		
		\textbf{Consequences of High False Positives:}
		\begin{itemize}
			\item Analyst fatigue - reviewers become desensitized
			\item Operational cost - each false positive costs $25-$75
			\item Legitimate customer friction - delayed payments, frozen accounts
			\item True positives may be missed in the noise
		\end{itemize}
		
		\examfocus{BTL testing is often overlooked by banks - exam tests whether you know it is required.}
	\end{frame}
	
	\begin{frame}{BTL Testing - Practical Example}
		\begin{block}{Example: Wire Transfer Scenario BTL Test}
			Testing a scenario designed to detect large wires to high-risk jurisdictions.
		\end{block}
		
		\textbf{Test Data Set:}
		\begin{itemize}
			\item 500 legitimate, low-risk customer transactions
			\item Payroll deposits, utility bills, mortgage payments
			\item All to low-risk jurisdictions (US, UK, Canada, EU)
			\item All amounts between $100-$5,000
		\end{itemize}
		
		\textbf{Scenario Being Tested:}
		\begin{itemize}
			\item Alert on any wire transfer > $2,500 to any jurisdiction
			\item No customer segmentation applied
		\end{itemize}
		
		\textbf{Results:}
		\begin{itemize}
			\item 85 of 500 transactions generated alerts (17 percent false positive rate)
			\item Legitimate wires > $2,500 to low-risk jurisdictions flagged
		\end{itemize}
		
		\textbf{Root Cause:}
		\begin{itemize}
			\item Scenario doesn't consider jurisdiction risk
			\item No segmentation by customer risk tier
			\item Threshold too low for legitimate business wires
		\end{itemize}
		
		\textbf{Remediation:}
		\begin{itemize}
			\item Apply tiered thresholds by jurisdiction risk
			\item High-risk: $1,000 threshold
			\item Medium-risk: $10,000 threshold
			\item Low-risk: $25,000 threshold
			\item Re-test: 2 percent false positive rate achieved
		\end{itemize}
		
		\redflag{BTL testing identified that jurisdiction risk must be incorporated into threshold setting.}
	\end{frame}
	
	\begin{frame}{ATL and BTL Testing - Combined Framework}
		\begin{block}{Balancing Detection and Efficiency}
			ATL and BTL testing together provide a complete picture of scenario performance.
		\end{block}
		
		\textbf{Performance Quadrants:}
		\begin{longtable}{p{0.25\textwidth}|p{0.35\textwidth}|p{0.35\textwidth}}
			\hline
			\rowcolor{lightgray}
			\textbf{Quadrant} & \textbf{ATL Result} & \textbf{BTL Result} \\
			\hline
			Ideal & High capture rate (90%+) & Low false positive rate (<10%) \\
			\hline
			Over-Filtered & Low capture rate (<90%) & Low false positive rate (<10%) \\
			\hline
			Under-Filtered & High capture rate (90%+) & High false positive rate (>10%) \\
			\hline
			Poor Design & Low capture rate (<90%) & High false positive rate (>10%) \\
			\hline
		\end{longtable}
		
		\textbf{Remediation by Quadrant:}
		\begin{itemize}
			\item Over-Filtered: Lower thresholds, broaden scenario logic
			\item Under-Filtered: Raise thresholds, narrow scenario logic, improve segmentation
			\item Poor Design: Complete scenario rewrite, reconsider typology approach
		\end{itemize}
		
		\textbf{Testing Frequency:}
		\begin{itemize}
			\item Full ATL/BTL testing: At least annually
			\item Trigger-based testing: After any material change to scenario, threshold, or typology
			\item Continuous monitoring: Track capture rate and false positive rate monthly
		\end{itemize}
		
		\examfocus{The exam tests that both ATL and BTL are required - not just one or the other.}
	\end{frame}
	
	\begin{frame}{Model Risk Management (MRM) Framework - SR 11-7}
		\begin{block}{Regulatory Requirements for Model Validation}
			All AML models (transaction monitoring, sanctions screening, fraud detection) are subject to model risk management under OCC 2011-12 and Federal Reserve SR 11-7.
		\end{block}
		
		\textbf{Key MRM Components:}
		\begin{enumerate}
			\item Model Inventory: All models documented, owners assigned, risk tier (high, medium, low)
			\item Model Development Documentation: Theoretical basis, assumptions, limitations, data sources
			\item Independent Validation: Before deployment and periodically after (annually for high-risk)
			\item Performance Monitoring: Ongoing tracking of precision, recall, false positive rate, drift
			\item Model Governance: Policies, procedures, approval authorities, escalation path
		\end{enumerate}
		
		\textbf{Validator Independence Requirements:}
		\begin{itemize}
			\item Validator cannot be involved in model development
			\item Validator must report functionally to independent risk management
			\item External validation for high-risk models (e.g., sanctions screening)
		\end{itemize}
		
		\redflag{A model that cannot be validated (black box AI) is a regulatory violation.}
	\end{frame}
	
	\begin{frame}{Model Risk Management - Validation Requirements}
		\begin{block}{What Model Validation Must Address}
			A complete validation report answers these specific questions.
		\end{block}
		
		\textbf{Validation Scope:}
		\begin{itemize}
			\item Conceptual Soundness: Does the model logic make sense for AML detection? Are assumptions reasonable?
			\item Data Quality: Is input data complete, accurate, and timely? Are there data gaps?
			\item Outcomes Analysis: What are precision, recall, ATL capture rate, BTL false positive rate?
			\item Benchmarking: How does this model compare to alternative approaches?
			\item Limitations and Compensating Controls: What can't the model detect? What other controls fill gaps?
		\end{itemize}
		
		\textbf{Validation Frequency by Risk Tier:}
		\begin{longtable}{p{0.25\textwidth}|p{0.35\textwidth}|p{0.35\textwidth}}
			\hline
			\rowcolor{lightgray}
			\textbf{Risk Tier} & \textbf{Full Validation} & \textbf{Ongoing Monitoring} \\
			\hline
			High Risk & Annually & Quarterly performance review \\
			\hline
			Medium Risk & Every 2 years & Semi-annual performance review \\
			\hline
			Low Risk & Every 3 years & Annual performance review \\
			\hline
		\end{longtable}
		
		\examfocus{Transaction monitoring models are typically classified as high or medium risk.}
	\end{frame}
	
	\begin{frame}{Model Risk Management - Validation Findings and Remediation}
		\begin{block}{Handling Validation Deficiencies}
			Finding a problem is not enough. There must be a documented remediation process.
		\end{block}
		
		\textbf{Severity Classifications:}
		\begin{itemize}
			\item Critical Finding: Material deficiency requiring immediate action. Model may need to be suspended.
			\item Major Finding: Significant deficiency with defined remediation timeline (30-90 days).
			\item Minor Finding: Improvement opportunity with longer remediation timeline (6-12 months).
			\item Observation: Suggestion for enhancement, no formal remediation required.
		\end{itemize}
		
		\textbf{Remediation Requirements:}
		\begin{itemize}
			\item Written remediation plan for each finding
			\item Assigned owner and due date
			\item Regular status updates to model governance committee
			\item Evidence of remediation (re-testing results)
			\item Closure approval by independent validator
		\end{itemize}
		
		\textbf{Common Validation Findings (Exam Scenarios):}
		\begin{itemize}
			\item Model drift - performance degraded without recalibration
			\item Data quality issues - incomplete or inaccurate input data
			\item Thresholds not calibrated - arbitrary values without ATL/BTL
			\item Segmentation inadequate - same rules for all customers
			\item Documentation missing - cannot explain model decisions
		\end{itemize}
		
		\redflag{Unremediated validation findings are a regulatory enforcement action waiting to happen.}
	\end{frame}
	
	\begin{frame}{Tuning Scenarios for Effectiveness - The Calibration Cycle}
		\begin{block}{Continuous Improvement Process}
			Scenario tuning is not a one-time event. It requires ongoing monitoring and adjustment.
		\end{block}
		
		\textbf{The Calibration Cycle:}
		\begin{enumerate}
			\item Monitor: Track alert volume, false positive rate, capture rate
			\item Analyze: Review sample of alerts, identify patterns
			\item Adjust: Modify thresholds, logic, or segmentation
			\item Validate: Run ATL/BTL testing on changes
			\item Deploy: Implement changes in production
			\item Repeat: Continue monitoring after deployment
		\end{enumerate}
		
		\textbf{Tuning Levers:}
		\begin{itemize}
			\item Threshold Values: Raise or lower monetary, velocity, or count thresholds
			\item Logic Operators: Add AND/OR conditions to narrow or broaden
			\item Time Windows: Expand or contract rolling periods
			\item Segmentation: Apply different rules to different customer types
			\item Exclusions: Remove known false positive patterns (e.g., internal transfers)
		\end{itemize}
		
		\examfocus{The exam tests that tuning must be documented and validated, not arbitrary.}
	\end{frame}
	
	\begin{frame}{Tuning Scenarios - False Positive Reduction}
		\begin{block}{Systematic Approach to Reducing False Positives}
			High false positive rates waste resources and cause analyst fatigue.
		\end{block}
		
		\textbf{False Positive Root Cause Analysis:}
		\begin{itemize}
			\item Review sample of 100-200 false positive alerts
			\item Categorize by root cause
			\item Common categories: customer type, transaction type, time of day, counterparty
			\item Identify patterns (e.g., payroll deposits always flagged)
		\end{itemize}
		
		\textbf{Tuning Solutions by Root Cause:}
		\begin{longtable}{p{0.35\textwidth}|p{0.6\textwidth}}
			\hline
			\rowcolor{lightgray}
			\textbf{Root Cause} & \textbf{Tuning Solution} \\
			\hline
			Payroll deposits flagged as structuring & Add exclusion logic for ACH credits with payroll description \\
			\hline
			Internal transfers between own accounts & Exclude same-owner accounts from velocity calculations \\
			\hline
			Business operating expenses flagged & Raise thresholds for business customers, add expense category exclusions \\
			\hline
			ATM withdrawals flagged as cash structuring & Separate ATM from teller cash, apply different thresholds \\
			\hline
			Low-risk jurisdiction wires flagged & Apply tiered thresholds by jurisdiction risk \\
			\hline
		\end{longtable}
		
		\textbf{Expected Outcome:}
		\begin{itemize}
			\item False positive rate reduction of 30-50 percent typical
			\item Maintain or improve capture rate (ATL)
			\item Document all exclusion logic for auditors
		\end{itemize}
		
		\redflag{Exclusion logic must be documented and validated - not just added arbitrarily.}
	\end{frame}
	
	\begin{frame}{Tuning Scenarios - Improving Capture Rate (Reducing False Negatives)}
		\begin{block}{Detecting What You're Missing}
			Low capture rate means criminals are operating undetected.
		\end{block}
		
		\textbf{Signs of Low Capture Rate:}
		\begin{itemize}
			\item ATL testing shows capture rate below 90 percent
			\item New SARs filed for typologies existing scenarios should detect
			\item Law enforcement inquiries about activity not alerted
			\item Peer institutions filing more SARs for similar typologies
		\end{itemize}
		
		\textbf{Tuning Solutions for Low Capture Rate:}
		\begin{itemize}
			\item Lower thresholds - but test impact on false positives first
			\item Broaden scenario logic - add more indicators
			\item Add new scenarios for emerging typologies
			\item Improve data quality - missing transactions or fields
			\item Update customer segmentation - ensure high-risk have lower thresholds
		\end{itemize}
		
		\textbf{Example: Structuring Scenario Tuning}
		\begin{itemize}
			\item Original: Alert on 8+ cash deposits < $10,000 in 30 days
			\item Capture rate: 82 percent (18 percent missed)
			\item After analysis: Criminals adapted to 5-6 deposits
			\item Tuned: Alert on 5+ cash deposits < $10,000 in 30 days
			\item New capture rate: 96 percent (4 percent missed)
			\item False positive impact: Acceptable increase of 15 percent
		\end{itemize}
		
		\examfocus{Sacrificing some false positive increase is acceptable to improve capture rate.}
	\end{frame}
	
	{frame}{Scenario Documentation Requirements}
	\begin{block}{What Regulators Expect to See for Each Scenario}
		Every monitoring scenario must have complete, up-to-date documentation.
	\end{block}
	
	\textbf{Required Documentation Elements:}
	\begin{itemize}
		\item Scenario Name and ID: Unique identifier for tracking
		\item Typology Detected: Which money laundering method this targets
		\item Risk Assessment Linkage: Reference to risk assessment entry
		\item Customer Segmentation: Which customer types this applies to
		\item Logic Statement: Exact rule in pseudo-code or SQL
		\item Threshold Values: All monetary, velocity, and count thresholds
		\item Time Windows: Rolling periods for aggregation
		\item Exclusions: Any transactions or customers excluded
		\item Expected Alert Volume: Monthly estimate by segment
		\item Validation History: ATL and BTL test results with dates
		\item Tuning History: All changes with dates and rationales
		\item Owner: Responsible analyst or team
	\end{itemize}
	
	\redflag{Missing or outdated scenario documentation is a common audit finding.}
	\end{frame}
	
	\begin{frame}{Alert Prioritization and Triage}
	\begin{block}{Not All Alerts Are Equal}
		Investigators must focus on highest-risk alerts first. Prioritization is critical.
	\end{block}
	
	\textbf{Prioritization Factors:}
	\begin{itemize}
		\item Customer Risk Tier: PEP > High > Medium > Low
		\item Transaction Amount: Larger amounts prioritized
		\item Jurisdiction Risk: High-risk countries prioritized
		\item Scenario Type: Certain typologies flagged as higher risk
		\item Velocity: Multiple alerts on same customer
		\item Recency: Newer alerts (FIFO or LIFO based on policy)
	\end{itemize}
	
	\textbf{Priority Scoring Example:}
	\begin{itemize}
		\item Score = (Customer Risk Weight x 5) + (Amount Score x 3) + (Jurisdiction Score x 2)
		\item High Priority (score > 80): Review within 24 hours
		\item Medium Priority (score 50-80): Review within 3 days
		\item Low Priority (score < 50): Review within 7 days
	\end{itemize}
	
	\textbf{Aging Requirements:}
	\begin{itemize}
		\item Regulators expect timely alert review
		\item Typical SLA: 90 percent of alerts reviewed within 30 days
		\item High-priority alerts: 100 percent within 48-72 hours
		\item Aging reports must be reviewed by management weekly
	\end{itemize}
	
	\redflag{Old, unreviewed alerts are a regulatory deficiency.}
	\end{frame}
	
	\begin{frame}{Alert Volume Management - Capacity Planning}
	\begin{block}{Matching Resources to Alert Volume}
		Alert volume must be aligned with investigator capacity.
	\end{block}
	
	\textbf{Capacity Calculation:}
	\begin{itemize}
		\item Average time per alert review: 15-30 minutes
		\item Average alerts per investigator per day: 16-32
		\item Effective working days per year: 220 (after training, meetings, leave)
		\item Annual capacity per investigator: 3,500-7,000 alerts
	\end{itemize}
	
	\textbf{Capacity Planning Formula:}
	\begin{itemize}
		\item Required Investigators = (Monthly Alert Volume x 12) / Annual Capacity per Investigator
		\item Example: 8,000 alerts/month x 12 = 96,000 alerts/year
		\item 96,000 / 5,000 (mid-range capacity) = 19.2 investigators
		\item Add 20% for management, training, quality assurance = 23 investigators
	\end{itemize}
	
	\textbf{Capacity Warning Signs:}
	\begin{itemize}
		\item Aging alerts exceed SLAs consistently
		\item Investigators working overtime regularly
		\item High turnover in monitoring team
		\item Quality review finds errors from rushing
		\item Management rejects requests for additional headcount
	\end{itemize}
	
	\redflag{Understaffing leading to unreviewed alerts is a regulatory violation.}
	\end{frame}
	
	\begin{frame}{Alert Disposition and Documentation}
	\begin{block}{Every Alert Must Have a Final Determination}
		No alert can be left unresolved. Each requires documented disposition.
	\end{block}
	
	\textbf{Valid Dispositions:}
	\begin{itemize}
		\item False Positive: No suspicious activity found. Document why.
		\item True Positive - No SAR: Suspicious but below filing threshold or mitigated. Document rationale.
		\item True Positive - SAR Filed: Suspicious activity meeting filing criteria. SAR submitted.
		\item Referred to Another Unit: e.g., Fraud, Legal, or Relationship Manager. Document referral.
		\item Pending Additional Information: Awaiting response from customer or internal source. Must have follow-up date.
	\end{itemize}
	
	\textbf{Investigation Documentation Requirements:}
	\begin{itemize}
		\item Customer information reviewed (KYC, CDD, EDD)
		\item Transaction details analyzed (amounts, dates, counterparties)
		\item External research performed (OSINT, databases)
		\item Analysis and rationale for disposition
		\item Supporting evidence attached (screenshots, reports)
		\item Supervisor approval for certain dispositions
		\item Audit trail of all actions and views
	\end{itemize}
	
	\redflag{Inadequate investigation documentation is a common quality assurance finding.}
	\end{frame}
	
	\begin{frame}{False Positive Analysis - Root Cause Categorization}
	\begin{block}{Systematic False Positive Reduction}
		Tracking false positive categories enables targeted tuning.
	\end{block}
	
	\textbf{Common False Positive Categories:}
	\begin{longtable}{p{0.3\textwidth}|p{0.35\textwidth}|p{0.3\textwidth}}
		\hline
		\rowcolor{lightgray}
		\textbf{Category} & \textbf{Example} & \textbf{Tuning Solution} \\
		\hline
		Payroll & Regular salary deposits flagged as structuring & Add payroll exclusion logic \\
		\hline
		Internal Transfers & Moving money between own accounts & Exclude same-owner accounts \\
		\hline
		Business Expenses & Rent, utilities, supplies flagged & Raise business thresholds \\
		\hline
		Low-Risk Jurisdiction & Wire to Canada flagged as high-risk & Tier thresholds by risk \\
		\hline
		Data Error & Corrupted field causes false match & Fix data quality issue \\
		\hline
		Threshold Too Low & $500 wire triggers alert & Raise threshold \\
		\hline
		Logic Too Broad & Any wire over $1,000 flagged & Narrow logic conditions \\
		\hline
		Segmentation Missing & Retail and corporate same rules & Implement segmentation \\
		\hline
	\end{longtable}
	
	\textbf{False Positive Tracking Metrics:}
	\begin{itemize}
		\item False positive rate by scenario (identify worst performers)
		\item False positive rate by customer segment (identify over-monitored segments)
		\item False positive rate by investigator (identify training needs)
		\item Root cause distribution (guide tuning priorities)
	\end{itemize}
	
	\examfocus{The exam tests that false positive analysis should drive scenario tuning.}
	\end{frame}
	
	\begin{frame}{Quality Assurance (QA) for Transaction Monitoring}
	\begin{block}{Ensuring Investigation Quality}
		Regular QA reviews validate that alerts are properly investigated and dispositions are correct.
	\end{block}
	
	\textbf{QA Sample Methodology:}
	\begin{itemize}
		\item Random sample of all closed alerts (typically 5-10 percent)
		\item Oversample high-risk alerts (100 percent of PEP alerts)
		\item Oversample SAR decisions (100 percent of SAR filings)
		\item Include false positives, true positives, and referrals
	\end{itemize}
	
	\textbf{QA Review Criteria:}
	\begin{itemize}
		\item Completeness: All required fields populated?
		\item Accuracy: Is disposition supported by evidence?
		\item Timeliness: Was alert reviewed within SLA?
		\item Documentation: Is rationale clear and complete?
		\item Consistency: Similar cases decided similarly?
	\end{itemize}
	
	\textbf{QA Findings Severity:}
	\begin{itemize}
		\item Critical: Missed SAR should have been filed. Immediate remediation.
		\item Major: Significant documentation gaps or procedural errors.
		\item Minor: Minor documentation issues, no impact on decision.
	\end{itemize}
	
	\redflag{QA program must be independent of the investigation team.}
	\end{frame}
	
	\begin{frame}{Back-Testing - Validating Historical Performance}
	\begin{block}{Learning from Past Mistakes}
		Back-testing runs historical transactions through current scenarios to identify missed suspicious activity.
	\end{block}
	
	\textbf{Back-Testing Methodology:}
	\begin{enumerate}
		\item Select historical period (e.g., previous 12 months)
		\item Run all transactions through current monitoring scenarios
		\item Identify transactions that would generate alerts today but did not at the time
		\item Review these missed transactions for suspicious activity
		\item File SARs for any suspicious activity identified
		\item Tune scenarios to prevent future misses
	\end{enumerate}
	
	\textbf{Trigger Events for Back-Testing:}
	\begin{itemize}
		\item New scenario deployment (test against historical data)
		\item Major threshold changes (validate improvement)
		\item After discovering a previously undetected typology
		\item Regulatory request or examination finding
		\item Material change to sanctions lists
	\end{itemize}
	
	\textbf{Regulatory Expectation:}
	\begin{itemize}
		\item Back-testing after significant scenario changes
		\item Documentation of review and any SAR filings
		\item No specific frequency mandated, but risk-based
	\end{itemize}
	
	\examfocus{Back-testing discovered the 18-month gap in the over-filtering trap scenario.}
	\end{frame}
	
	\begin{frame}{Look-Back Reviews - Regulatory Requirements}
	\begin{block}{When Regulators Require Historical Reviews}
		Look-back reviews examine historical transactions after discovering a control failure.
	\end{block}
	
	\textbf{Common Look-Back Triggers:}
	\begin{itemize}
		\item Silent data feed failure (transactions never monitored)
		\item Scenario logic error (wrong threshold or condition)
		\item Missed sanctions list update (screening using old list)
		\item Model drift without recalibration
		\item Regulatory finding requiring historical review
	\end{itemize}
	
	\textbf{Look-Back Scope Determination:}
	\begin{itemize}
		\item Time period of control failure (must be identified)
		\item Transaction types affected (wires, cash, ACH, etc.)
		\item Customer segments affected (may be all or subset)
		\item Regulator may mandate specific scope
	\end{itemize}
	
	\textbf{Look-Back Execution:}
	\begin{itemize}
		\item Obtain all missed transactions (may require IT extraction)
		\item Run through current monitoring scenarios
		\item Investigate all alerts generated
		\item File SARs for any suspicious activity identified
		\item Document all findings for regulators
	\end{itemize}
	
	\redflag{Failure to conduct a look-back after discovering a control failure is a separate violation.}
	\end{frame}
	
	\begin{frame}{Key Performance Indicators (KPIs) for Transaction Monitoring}
	\begin{block}{Metrics Management Must Review Monthly}
		Regular KPI reporting enables oversight of monitoring effectiveness.
	\end{block}
	
	\textbf{Volume Metrics:}
	\begin{itemize}
		\item Total transactions monitored (by channel, product, segment)
		\item Total alerts generated (by scenario, priority, segment)
		\item Alert rate (alerts per 1,000 transactions)
		\item SARs filed (by typology, segment, scenario)
		\item SAR filing rate (SARs per 1,000 alerts)
	\end{itemize}
	
	\textbf{Performance Metrics:}
	\begin{itemize}
		\item False positive rate (by scenario, segment)
		\item Capture rate (from ATL testing)
		\item Alert aging (percentage reviewed within SLA)
		\item Investigation time (average hours per alert)
		\item Quality assurance scores
	\end{itemize}
	
	\textbf{Resource Metrics:}
	\begin{itemize}
		\item Alerts per investigator
		\item Cost per alert reviewed
		\item Investigator turnover rate
		\item Training hours completed
		\item Overtime hours
	\end{itemize}
	
	\textbf{Reporting Cadence:}
	\begin{itemize}
		\item Daily: Alert volume, aging, priority distribution
		\item Weekly: Investigator workload, emerging patterns
		\item Monthly: Complete KPI dashboard to management
		\item Quarterly: Trend analysis to Board
	\end{itemize}
	
	\examfocus{Management must receive and act on KPI reports. Ignoring negative trends is a deficiency.}
	\end{frame}
	
	\begin{frame}{Common Regulatory Findings in Transaction Monitoring}
	\begin{block}{What Regulators Cite in Enforcement Actions}
		These deficiencies appear repeatedly in consent orders.
	\end{block}
	
	\textbf{Most Common Findings:}
	\begin{itemize}
		\item Over-Filtering: Thresholds set too high, causing false negatives
		\item Data Quality: Incomplete or inaccurate data feeds
		\item Scenario Gaps: Missing scenarios for relevant typologies
		\item No ATL/BTL Testing: Cannot demonstrate scenario effectiveness
		\item Outdated Scenarios: Not updated for emerging typologies
		\item Inadequate Segmentation: Same rules for all customer types
		\item Alert Backlog: Unreviewed alerts exceeding SLAs
		\item Poor Documentation: Missing rationale for dispositions
		\item No Independent Validation: Models deployed without MRM
		\item Model Drift: Performance degraded without recalibration
	\end{itemize}
	
	\textbf{Enforcement Consequences:}
	\begin{itemize}
		\item Civil money penalties (millions to billions)
		\item Consent orders with mandated remediation
		\item Independent monitor required
		\item Board and management reporting requirements
		\item Restrictions on growth or acquisitions
	\end{itemize}
	
	\redflag{Most enforcement actions include at least three of these findings.}
	\end{frame}
	
	\begin{frame}{Transaction Monitoring Program - Minimum Requirements Checklist}
	\begin{block}{Regulatory Expectations Summary}
		Every transaction monitoring program must include these elements.
	\end{block}
	
	\textbf{Required Program Elements:}
	\begin{enumerate}
		\item Written policies and procedures for monitoring
		\item Risk assessment identifying relevant typologies
		\item Scenarios mapped to each relevant typology
		\item Documented thresholds with calibration methodology
		\item Customer segmentation logic
		\item ATL and BTL testing (at least annually)
		\item Model risk management and independent validation
		\item Alert prioritization and triage process
		\item Investigation procedures and documentation standards
		\item Quality assurance program
		\item Management reporting and oversight
		\item Regular scenario tuning and updates
		\item Data quality controls and reconciliation
		\item Audit trail of all monitoring activities
		\item Back-testing and look-back procedures
	\end{enumerate}
	
	\examfocus{The exam tests that all these elements are required, not optional.}
	\end{frame}
	
	\begin{frame}{Transaction Monitoring - Exam Day Summary}
	\begin{itemize}
		\item Transaction monitoring is the primary detective control for ongoing AML
		\item Rules-based systems are explainable but miss new typologies
		\item Customer segmentation prevents one-size-fits-all monitoring
		\item Thresholds must be data-driven, not arbitrary
		\item ATL tests effectiveness (capture rate)
		\item BTL tests efficiency (false positive rate)
		\item Both ATL and BTL are required - not just one
		\item Over-filtering (thresholds too high) creates false negatives
		\item False negatives are worse than false positives
		\item Model risk management requires independent validation
		\item Model drift requires regular recalibration
		\item Scenarios must be tuned continuously based on performance data
		\item Documentation is as important as detection
		\item Alert backlogs are a regulatory violation
		\item Quality assurance validates investigation quality
	\end{itemize}
	
	\greennote{Master these concepts for the 6-8 transaction monitoring questions on the exam.}
	\end{frame}
	
\end{document}